\documentclass[pdftex,12pt]{article}
\usepackage{latexsym,fancyhdr}
\usepackage[a4paper,includeheadfoot,left=1.8cm,right=1.2cm,,top=2cm,bottom=2cm,twoside]{geometry}
\usepackage[lyric]{songs}
\usepackage[utf8]{inputenc}
\DeclareUnicodeCharacter{FEFF}{}
\usepackage[english,danish]{babel}
%\usepackage[english]{babel}
\usepackage{amssymb}
\usepackage{stmaryrd}
\usepackage{amsmath}
\usepackage{amsthm}
\usepackage[final]{pdfpages}
\usepackage[none]{hyphenat}
\usepackage{hyperref}
\usepackage{multicol}
\input binhex
\setlength{\columnsep}{1.5cm}
\newindex{titleidx}{titlefile}
\sepindexesfalse
\noversenumbers
\title{Sangbog Matematik \the\year}
\author{Ungdommens Naturvidenskabelige Forening}
\date{\today}
%\addtolength{\headwidth}{\marginparsep}
%\addtolength{\headwidth}{\marginparwidth}
\renewcommand{\headrulewidth}{.4pt}
%\renewcommand{\footrulewidth}{0.4pt}
\fancyhead{}
\fancyhead[CE,CO]{}
\fancyhead[RE,LO]{\thepage}
\fancyfoot{}
\pagestyle{fancy}
\newcommand{\countstyle}{greek}
\makeatletter

\newcommand*{\binaryX}[1]{%
  \expandafter\@binaryX\csname c@#1\endcsname
}

\newcommand*{\@binaryX}[1]{0b\binary{#1}}

\newcommand*{\octX}[1]{%
  \expandafter\@octX\csname c@#1\endcsname
}

\newcommand*{\@octX}[1]{0o\oct{#1}}

\newcommand*{\hexX}[1]{%
  \expandafter\@hexX\csname c@#1\endcsname
}

\newcommand*{\@hexX}[1]{0x\hex{#1}}

\newcommand*{\greek}[1]{%
  \expandafter\@greek\csname c@#1\endcsname
}

\newcommand*{\@greek}[1]{%
  $\ifcase#1
  %\or
  \alpha\or\beta\or\gamma\or\delta\or\varepsilon
  \or\zeta\or\eta\or\theta\or\iota\or\kappa\or\lambda
  \or\mu\or\nu\or\xi\or o\or\pi\or\rho\or\sigma
  \or\tau\or\upsilon\or\varphi\or\chi\or\psi\or\omega
  \or A\or B\or\Gamma\or\Delta\or E
  \or\Zeta\or H\or\Theta\or I\or K\or\Lambda
  \or M\or N\or\Xi\or O\or\Pi\or P\or\Sigma
  \or T\or\Upsilon\or\Phi\or X\or\Psi\or\Omega
  \else\@ctrerr\fi$
}


\newcommand*{\symcount}[1]{%
  \expandafter\@symcount\csname c@#1\endcsname
}

\newcommand*{\@symcount}[1]{%
  $\ifcase#1
  \neg\or\sqrt{\ }\or \int
  \or +\or \times\or\pm\or \circ\or \oplus\or \otimes
  \or\cup \or\cap%\or\sqcup\or\sqcap
  \or \vert\or \Vert\or \perp%\or\angle
  \or\sim
  \or <\or >
  \or =\or \cong
  \or \in
  \or\exists\or\forall\or \vee\or \wedge
  \or\Rightarrow\or\mapsto\or\Leftrightarrow
  \or\lightning\or\square
  \else\@ctrerr\fi$
}

\makeatother

\pagenumbering{shiftedpage}
\begin{document}
\newcounter{temp}
\newcounter{temppage}
\newcounter{pageoffset}
\setcounter{page}{0}
\newgeometry{margin=.5cm,top=4cm,bottom=.5cm}
\thispagestyle{empty}
\addtocounter{pageoffset}{-1}
\centering
\phantom{test}

\mbox{\includegraphics[scale=0.557376932846]{Unflogo.pdf}}
\vspace{.5cm}
\begin{center}
\fontfamily{phv}\fontsize{50}{60}\selectfont Matematik Camp\\\the\year

\vspace{1cm}
\fontsize{35}{40}\selectfont Sangbog\end{center}
\restoregeometry
\raggedright
\songpos{0}
\spenalty=50
\vvpenalty=100
\vcpenalty=100
\cvpenalty=100
\ccpenalty=100
\interlinepenalty=200
\setlength{\parindent}{1cm}
\begin{songs}{}
\setcounter{songnum}{0}

\beginsong{Vi kan ikke li'}%0
[sr={Melodi: Vores buschauffør kan ikke køre bus}
,
by={}
,
cr={}]\hypertarget{Vi kan ikke li'}{}
\label{song0}

\beginverse
Vi kan ikke li' folk vi ik' kan li'
Vi kan ikke li' folk vi ik' kan li'
For vi kan sgu ikke li' dem,
de skulle hellere tage at blive hjemme
Vi kan ikke li' folk vi ik' kan li'

\endverse
\beginverse
Vi kan ikke li' folk fra teologi\ldots
For de er nogle hængemuler,
tror at helligånden puler

\endverse
\beginverse
Vi kan ikke li' folk fra biologi\ldots
For de går i strikketrøjer,
og med dyrene sig fornøjer

\endverse
\beginverse
Vi kan ikke li' folk der har fysik\ldots
For de har sgu nogle testikler
der er på størrelse med partikler

\endverse
\beginverse
Vi kan ikke li' folk fra musikvidenskab\ldots
For de er nogle lamme nødder,
passer ikke på versefødder

\endverse
\beginverse
Vi kan ikke li' folk fra nano-tek\ldots
De går rundt med sure miner,
laver bittesmå maskiner

\endverse
\beginverse
Vi kan ikke li' folk fra statistik\ldots
De vil hypoteser teste
mens de hel're burde feste

\endverse
\beginverse
Vi kan ikke li' folk fra arkæologi\ldots
fordi de graver og de graver
og på gamle mennesker rager

\endverse
\beginverse
Vi kan ikke li' folk fra filosofi\ldots
for de går bar' og funderer
om de selv mon eksisterer

\endverse
\beginverse
Vi kan ikke li' folk fra farmaci\ldots
for de går og triller piller
og tabletter og pastiller

\endverse
\beginverse
Vi kan ikke li' folk fra geologi\ldots
For de går og ser på stenene
når de burde sprede benene

\endverse
\beginverse
Vi kan ikke li' folk fra astronomi\ldots
De er fjolser uden hjerner
der går rundt og ser på stjerner

\endverse
%\beginverse
%Vi kan ikke li' folk fra iNano\ldots
%bagfra blir' de lidt for kække
%synes Ångstrøm de er frække

%\endverse
\beginverse
Vi kan ikke li' folk fra økonomi\ldots
For de regner på budgettet
når de sidder på toilettet

\endverse
\beginverse
Vi kan ikke li' folk fra lingvistik\ldots
Kulli waffli zarka gunku
emfle birnan smöja dunku

\endverse
\beginverse
Vi kan ikke li' folk der har kemi\ldots
For de hedder Ion og Ester
mens de luften stærkt forpester

\endverse
\beginverse
Vi kan ikke li' folk fra matematik\ldots
For de teoretiserer
indtil $2$ plus $3$ bli'r $4$

\endverse
\beginverse
Vi kan ikke li' folk der har IT\ldots
Uni skaffer studiejobbet
Det er næsten alt for snobbet

\endverse
\beginverse
Vi kan ikke li' folk der har dat-mult\ldots
Datalogcirkusartister;
de er næsten humanister

\endverse
%\beginverse
%Vi kan ikke li' folk der har tek-fys\ldots
%for de har en bærbar PC
%læser mails når de' på WC

%\endverse
%\beginverse
%Vi kan ikke li' folk fra statskundskab\ldots
%for de gider kun at kom' for- di
%de tror de møder Lomborg

%\endverse
%\beginverse
%Vi kan ikke li' folk fra farmaceutisk kemi\ldots
%de fungerer som en buffer
%der i alle emner skuffer

%\endverse
\beginverse
Vi kan ikke li' folk fra medicin\ldots
For de går i hvide kitler,
og det rimer jo på Hitler

\endverse
\beginverse
Vi kan ikke li' folk fra molekylær(biologi)\ldots
for de kloner små kaniner
splejser dem med appelsiner

\endverse
\beginverse
Vi kan ikke li' folk fra psykologi\ldots
For de vader rundt i læder,
og de hader deres fædre

\endverse
\beginverse
Vi kan ikke li' folk der er jurister\ldots
for med hule paragraffer
dagen lang sig selv de straffer

\endverse
\beginverse
Vi kan ikke li' folk der har historie\ldots
for de har kun gamle minder
og et job de aldrig finder

\endverse
\beginverse
Vi kan ikke li' folk der har idræt\ldots
For de ligner jo rødbeder,
når de render rundt og sveder

\endverse
\beginverse
Vi kan ikke li' folk fra datalogi\ldots

\endverse
\endsong


\beginsong{Den kanoniske matematikersang}%1
[sr={Melodi: Bamses fødselsdag}
,
by={Esben "Med Hår" Bistrup Halvorsen og Rasmus "Rus" Resen Amossen}
,
cr={}]\hypertarget{Den kanoniske matematikersang}{}
\label{song1}

\beginverse
Som mat'matikstuderende, så er jeg svær at narre.
Når andre kalder mig for nørd, så må jeg bare svare:
"Hip hurra for algebra,
Euler, Gauss og Galois
og for rum med en kompakt deformationsretrakt".

\endverse
\beginverse
En dag jeg sa' til kæresten "nu vil jeg dyrke grupper".
Hun kiggede forbavset op, "det ikke super duper!".
"Jamen det er ej med dig,
gutterne de er på vej".
"Får du ikke nok af mig?!" -så rejste hun sin vej.

\endverse
\beginverse
Nu var jeg blevet singelton og sagde til min moder:
"Jeg er disjunkt med kæresten som Peking med Nyboder.
Inklusionen er nu vendt,
jeg er blevet transcendent".
"Er du trans, din klamme tøs?!" -så blev jeg arveløs.

\endverse
\beginverse
Da arven nu var faldet bort, jeg måtte til at spare.
Jeg talte med min vicevært og kunne ham forklare:
"Pengemængden er kompakt,
jeg vil ha' en ny kontrakt!".
"Fint med mig, den kommer her: Du bor her ikke mer'!".

\endverse
\beginverse
Nu var jeg efterhånden ved at få lidt hovedpine.
Jeg tog til hospitalet og fik hjælp af en blondine.
"Hovedsmerten går for vidt,
den sku' deles op i snit!".
"Hvidt og snit, og så nå'et, ik'?" -der røg mit overblik.

\endverse
\beginverse
Jeg har nu få't det hvide snit og boligen er røget,
men selvom både arv og kær'ste fløj, er jeg fornøjet!
Hip hurra for algebra,
Euler, Gauss og Galois
og for rum med en kompakt deformationsretrakt.

\endverse
\endsong



\beginsong{Jeg er en matematiker fra HCØ}%2
[sr={Melodi: Jeg er en papegøje fra Amerika}
,
by={Kennie Nybo Mortensen, Jens Gustav Jensen, Kennet Andersen, Christian C. Harhoff og Povl Aage}
,
cr={}]\hypertarget{Jeg er en matematiker fra HCØ}{}
\label{song2}

\beginverse
Jeg er en matematiker fra HCØ;
jeg var en tur i byen det var dejligt.
Da mødte jeg en pige, hun var rigtig sød,
men hør nu her hvordan det blev forfejlet.
Hun spurgte mig, åh falderi åh faldera,
``hva' laver du når du ikke er på bar?''
Da svared' jeg, åh falderi åh faldera,
``jeg skriver mit speciale i $C^*$-algebra.''

\endverse
\beginverse
``Hvor interessant med stjerner er du astrolog?
Hvad sker der næste måned, jeg er tvilling?''
``Det rimer fint, for jeg er nemlig pædagog.''
Her brød jeg ind, fortalte om min stilling.
Så sagde jeg, åh falderi åh faldera,
``du misforstod mig,'' for jeg var jo rar;
og fortsatte, åh falderi åh faldera,
``det drejer sig om operatoralgebra.''

\endverse
\beginverse
``Aha, hva' si’r du algebrator opera?
Så må du være noget ved musikken.
Men stjerner og musik, hvad er nu det, forklar!''
Da måtte jeg dosere matematiken.
Så sagde jeg, åh falderi åh faldera,
``tag $B(H)$, så er vi næsten klar,
hvor $H$ det er, åh falderi åh faldera,
et Banachrum med prikprodukt, halvandenlinear.''

\endverse
\beginverse
``BH'en af i Bjarnes rum, og Bjarne hvem?!
Er Bjarne en af dine sjofle venner?
Næ, jeg må hel're se at vende snuden hjem;
ta'r ikke hjem til en jeg ikke kender.''
Så råbte hun, åh falderi åh faldera,
``er du tenor, eller er jeg blot til nar?!''
Så tænkte jeg, åh falderi åh faldera,
``hvis det skal blive til noget, må beskeden være klar.''

\endverse
\beginverse
Så sagde jeg, og det var løgn, ``jeg synger bas
i komponisten Eulers sidste stykke.''
Hun smilte sødt. ``Det ærlig snak, ej mere gas.
Tag med mig hjem, der vil du gøre lykke!''
Moralen er, åh falderi åh faldera:
Hvis du skal score når du er på bar,
så tal ej sandt, åh falderi åh faldera.
Fortæl hende hel're at du spiller elguitar.

\endverse
\endsong



\beginsong{Hymne til Matematik Camp}%3
[sr={Melodi: I Danmark er jeg født}
,
by={Søren Møller, Matematik Camp 2014}
,
cr={}]\hypertarget{Hymne til Matematik Camp}{}
\label{song3}

\beginverse
På campen har jeg lært, der har jeg hjemme
Der er alt sandt, derfra min fremtid går.
Du mat'matik, du er min egen stemme,
så klart beviserne min hjerne når.
Min tankes $\pi$ og $\rho$,
hvor Euklids elementer
står mellem algebra og rentesrenter.
Dig agter jeg -- MatCamp mit origo.
\endverse

\beginverse
Hvor lægges næste år mon sommerskolen?
Mon ik' det er ved jydens Unisø.
Et yndigt sted hvor alle knækker koden,
så dejligt som på Danas egen ø.
Min tankes $\pi$ og $\rho$,
hvor aksiomer holder.
Gauss gav os alt, han er sandhedens tolder.
Dig agter jeg -- MatCamp mit origo.
\endverse

\beginverse
Engang du herre var i kloges færden.
Bød over dumhed - nu du kaldes led.
Et lille fag som dog ved alt om verden,
og hører vores sang og troskabsed.
Min tankes $\pi$ og $\rho$
funktioners nulpunkt finder.
Gauss giv dig rødder, som vi gi'r dig minder.
Dig agter jeg -- MatCamp mit origo.
\endverse

\beginverse
Se divergens ved polyedrets kanter,
der gi'r uend'ligt mange mindre skridt;
og symmetri blandt vores kommutanter
gør stadig majoranten lige skidt.
Min tankes $\pi$ og $\rho$,
hvor grænser frit bestemmes,
og ingen sig ved udvælgelsen græmmes.
Dig agter jeg -- MatCamp mit origo
\endverse

\beginverse
Du camp hvor jeg har lært, hvor jeg har hjemme,
hvor alt er sandt, hvorfra min fremtid går,
hvor mat'matik er livets klare stemme,
og klart beviserne min hjerne når
Min tankes $\pi$ og $\rho$
ved videnskabens kilde.
I mat'matik min hjerne alt vil vide.
Dig agter jeg -- MatCamp mit origo.
\endverse
\endsong



\beginsong{Når Jeg Kloner Min Kat}[sr={Melodi: Westlife - When You're Looking Like That}
,
by={BIOrevy 2012}
,
cr={}]\hypertarget{Når Jeg Kloner Min Kat}{}
\label{song4}

\beginverse
Når jeg kloner min kat

\endverse
\beginverse
Med en ny teknik kan os på bio skabe liv
Den ku' bruges ondt og klone Henrik Busch
Til vi har en hær
Det er ikke det, vi gør!
Bruger teknikken til at redde arter, men
Sker der noget skidt med dyret vi elsker, så

\endverse
\beginverse
Må mammutten vente
Ikke tid til at grundforske

\endverse
\beginverse
Ska' aldirg mere sig' farvel
Når jeg kloner min kat
Vil altid være perfekt, nuttet, sød
For altid vær' min skat
Aldrig købe mig en ny
Aldrig købe et erstatningskæledyr
Ska' aldrig mere sig farvel
Når jeg kloner min kat

\endverse
\beginverse
Jurassic Park var tyd'ligvis en succes
Dinosaurus på jord ku' være fedt
Men sig mig, hvordan gør man det?
Indsæt gen og skab superarter
Nem vej til at skabe über seje dyr
Men når katten forsvinder, har vi et hyr

\endverse
\beginverse
Så mammutten venter
ikke tid til at grundforske

\endverse
\beginverse
Ska' aldirg mere sig' farvel
Når jeg kloner min kat
Vil altid være perfekt, nuttet, sød
For altid vær' min skat
Aldrig købe mig en ny
Aldrig købet et erstatningskæledyr
Ska' aldrig mere sig farvel
Når jeg kloner min kat

\endverse
\beginverse
Skal jeg så spørge jer nu
Hvilket studie har magt?
Biologer har livets kraft
Så jeg kloner min kat

\endverse
\beginverse
Aldirg mere sig' farvel
Når jeg kloner min kat
Vil altid være perfekt, nuttet, sød
For altid vær' min skat
Aldrig købe mig en ny
Aldrig købet et erstatningskæledyr
Ska' aldrig mere sig farvel
Når jeg kloner min kat

\endverse
\beginverse
Ska' aldirg mere sig' farvel
Når jeg kloner min kat
Vil altid være perfekt, nuttet, sød
For altid vær' min skat
Aldrig købe mig en ny
Aldrig købet et erstatningskæledyr
Ska' aldrig mere sig farvel
Når jeg kloner min kat

\endverse
\beginverse
Ska' aldirg mere sig' farvel
Når jeg kloner min kat
Vil altid være perfekt, nuttet, sød
For altid vær' min skat
Aldrig købe mig en ny
Aldrig købet et erstatningskæledyr
Ska' aldrig mere sig farvel
Når jeg kloner min kat

\endverse
\endsong



\beginsong{Forever kvant}[sr={Melodi: Forever Young}
,
by={FysikRevy 2011}
,
cr={}]\hypertarget{Forever kvant}{}
\label{song5}

\beginverse
Jeg er til revy, det er min første gang,
tiden til Kvant 1 den virker pludselig lang.
Jeg griner med for jeg synes det' flovt,
at jeg ik' forstår at kvant er sjovt.
\endverse

\beginverse
Jeg skråler med på ``Kvanter i måneskin'',
men egentlig så fatter jeg ik' en pind.
Ja, jeg kan hverken regne ud eller ind
på energi og kvantespin.
\endverse

\beginverse
Tænk at forstå blot en enkel sang,
grine af en sketch blot en enkel gang,
vide hvem Per Hedegård egentlig er.
De andre griner,
men jeg
\endverse

\beginverse
forstår det ik',
før jeg får kvantemekanik!
Alle andre griner højt omkring mig,
men ik' mig,
men ik' mig!
\endverse

\beginverse
Forstår det ik',
før jeg får kvantemekanik!
Hva' en hæveoperator egentlig?
Jeg ved det ik'.
\endverse

\beginverse
Newtons tre love kan jeg godt forstå,
og termodynamik ku' jeg bestå.
Jeg har før hørt om Schrödinger, men
hvem er Hr. Milton? (Hamilton)
\endverse

\beginverse
Sku' lære LinAlg, men ved ikke hvorfor.
Er deltafunktion noget jeg får brug for?
Hvordan kan tid være imaginær?
Og hvor er det der Hilbertrum?
\endverse

\beginverse
Kan bra og ket hjælpe med at projicer'?
Hvad sker der når man pertuberer mer'?
Hvor'n er det at kvante i måneskin?
Er jeg til grin?
\endverse

\beginverse
Forstår det ik',
før jeg får kvantemekanik!
Alle andre griner højt omkring mig,
men ik' mig,
men ik' mig!
\endverse

\beginverse
Forstår det ik',
før jeg får kvantemekanik!
Alle andre griner højt omkring mig,
men ik' mig,
men ik' mig!
Forstår det ik',
for jeg har ik' haft nok fysik.
Ikke nok til FysikRevyen.
\endverse
\endsong



\beginsong{Afrundingsfejl}[sr={Melodi: Shakira - Whenever, Wherever}
,
by={DIKUrevy 2015}
,
cr={}]\hypertarget{Afrundingsfejl}{}
\label{song6}

\beginverse
        Går på nettet for at købe mig et
        nyt og lækkert gaming headset
        Finder et i en online forhandler
        Det var rim¿lig billigt -- jeg ta'r det

\endverse
\beginverse
        Trykker på check ud for at betale
        Der ses det frygteligt unormale
        Prisen vist med 8 decimaler
        Der er da vist brugt brudne tal der


\endverse
\beginverse
        Nej, nej, nej, nej, nej, nej
        Nej, nej, nej, nej, nej, nej
        Det gi'r jo\ldots
        Afrundingsfe-ejl!


\endverse
\beginverse
        Lidt over, lidt under
        Når du dit tal afrunder
        Du må ikke dividere
        Brug dog et heltal hér!

\endverse
\beginverse
        Produkter dig tugter
        Din præcision fordufter
        Tro nu ikke vi ej ser
        De floating point værdier


\endverse
\beginverse
        0,99997, 0,99997, 0,99997, 0,99997


\endverse
\beginverse
        Laver mad på DIKU, lægger ud så
        Jeg siger folk kan overføre
        Hører pengene er sendt og venter
        Logger ind på webbank og checker

\endverse
\beginverse
        Trykker mig nu ind på oversigten
        Ak, der begynder tillidssvigten
        Selvom beløbet er et heltal
        Er der fejl på ottende decimaltal


\endverse
\beginverse
        Nej, nej, nej, nej, nej, nej
        Nej, nej, nej, nej, nej, nej
        Der er jo\ldots
        Afrundingsfe-ejl!


\endverse
\beginverse
        Lidt over, lidt under
        Når du dit tal afrunder
        Du må ikke dividere
        Brug dog et heltal hér!

\endverse
\beginverse
        Produkter dig tugter
        Din præcision fordufter
        Tro nu ikke vi ej ser
        De floating point værdier


        0,99997, 0,99997, 0,99997, 0,99997


\endverse
\beginverse
        Nej, nej, nej, nej, nej, nej
        Nej, nej, nej, nej, nej, nej
        Tænk dig om
        Brug din fornu-uft

\endverse
\beginverse
        Nej, nej, nej, nej, nej, nej, nej, nej
        Pris bør gemmes i
        Typer med
        Fast præcisio-on!


\endverse
\beginverse
        Lidt over, lidt under
        Når du dit tal afrunder
        Du må ikke dividere
        Brug dog et heltal hér!

\endverse
\beginverse
        Produkter dig tugter
        Fejlen den akkumulér'
        Tro ikke vi ej ser
        Dine floating point værdier


\endverse
\beginverse
        Lidt over, lidt under
        Når du dit tal afrunder
        Du må ikke dividere
        Brug dog et heltal hér!

\endverse
\beginverse
        Produkter dig tugter
        Fejlen den akkumulér'
        Tro ikke vi ej ser
        Dine floating point værdier

\endverse
\endsong



\beginsong{Midsommervisen -- Vi elsker vort land}[sr={}
,
by={Holger Drachmann}
,
cr={}]\hypertarget{Midsommervisen -- Vi elsker vort land}{}
\label{song7}

\beginverse
Vi elsker vort land,
når den signede jul
tænder stjernen i træet med glans i hvert øje,
når om våren hver fugl
over mark, under strand
lader stemmen til hilsende triller sig bøje:
vi synger din lov over vej, over gade,
vi kranser dit navn, når vor høst er i lade;
\lrep men den skønneste krans
bli’r dog din, Sankte Hans!
den er bundet af sommerens hjerter så varme, så glade.\rrep
\endverse

\beginverse
Vi elsker vort land,
men ved midsommer mest,
når hver sky over marken velsignelser sender,
når af blomster er flest,
og når kvæget i spand
giver rigeligst gave til flittige hænder!
når ikke vi pløjer og harver og tromler,
når koen sin middag i kløveren gumler:
\lrep da gå ungdom til dans
på dit bud, Sankte Hans!
ret som føllet og lammet, der frit over engen sig tumler.\rrep
\endverse

\beginverse
Vi elsker vort land,
og med sværdet i hånd
skal hver udenvælts fjende beredte os kende,
men mod ufredens ånd
over mark, under strand
vil vi bålet på fædrenes gravhøje tænde:
Hver by har sin heks og hvert sogn sine trolde,
dem vil vi fra livet ved glædesblus holde;
\lrep vi vil fred hertillands,
Sankte Hans, Sankte Hans!
den kan vindes, hvor hjerterne aldrig bli’r tvivlende kolde!\rrep
\endverse
\endsong



\beginsong{Få det læst}[sr={Melodi: The Beatles - Let it Be}
,
by={Matematikrevyen 2006}
,
cr={}]\hypertarget{Få det læst}{}
\label{song8}

\beginverse
Når jeg sidder og skal læse lektier bliver jeg en doven rad
Jeg kan ikke ta’ det, få det læst!
Jeg bliver så let lidt distraheret, tv, spil og aftensmad
Ta’ dig nu lidt sammen, få det læst!
\endverse

\beginverse
Få det læst, få det læst, få det læst, få det læst
Ta’ dig nu lidt sammen, få det læst.
\endverse

\beginverse
Jeg syn’s at matematik er herligt, frem for alle andre fag
Så jeg bør jo kunne få det læst.
Men der er andet man kan lave, drikke løs og fyr’ den af
Jeg kan ikke klare det, få det læst!
\endverse

\beginverse
Få det læst, få det læst, få det læst, få det læst!
Jeg kan ikke klare det, få det læst.
Få det læst, få det læst, få det læst, få det læst!
Ta’ dig nu lidt sammen, få det læst.
\endverse

\beginverse
Og selv når jeg har læseferie, går det ikke særlig nemt
Venter til det sidste, få det læst!
Men jeg er tvunget for at huske tusind ting som jeg har glemt
Se nu på de formler, få det læst!
\endverse

\beginverse
Få det læst, få det læst, få det læst, få det læst!
Se nu på de formler, få det læst!
Få det læst, få det læst, få det læst, få det læst!
Ta’ dig nu lidt sammen, få det læst!
Få det læst, få det læst, få det læst, få det læst!
Hvis du vil bestå så få det læst!
\endverse
\endsong



\beginsong{Jeg vil læse *}[sr={Melodi:Regnvejrsdag i november}
,
by={DIKUrevy 2008}
,
cr={}]\hypertarget{Jeg vil læse *}{}
\label{song9}

\beginverse
Jeg vil læse retorik
for jeg fatter ik' en brik
jeg vil lær' om sprogfinesser
og gå op i petitesser
jeg bli'r nok en stor professor
jeg vil læse retorik

\endverse
\beginverse
Jeg vil være teolog
selvom lønnen ik' er god
tro på gud det er det sande
også når man træder vande
missionére i fjerne lande
jeg vil være teolog

\endverse
\beginverse
Jeg vil læse medicin
tjene penge som et svin
jeg vil skær i døde men'sker
men jeg er jo også svensker
klæde mig i \LaTeX handsker
jeg vil læse medicin

\endverse
\beginverse
Jeg vil læse på kemi
hårde stoffer kan jeg li'
jeg er skæv fra ni til sytten
finsprit hælder jeg i bøtten
så jeg ender nok på støtten
jeg vil læse på kemi

\endverse
\beginverse
Jeg vil læse cand.polit
for mit men'skesyn er skidt
jeg vil skær på hospitaler
og på skoler uden kvaler
mens jeg holder falske taler
jeg vil læse cand.polit

\endverse
\beginverse
Jeg vil læse på fysik
lære kvantemekanik
jeg går op Einsteins tanker
og jeg er en rigtig dranker
kæler for mit fadølsanker
Jeg vil læse på fysik

\endverse
\beginverse
Jeg vil være officér
i den danske dronnings hær
jeg vil skyde talibaner
vifte med de danske faner
sprede død i lange baner
jeg vil være officér

\endverse
\beginverse
Historie er mit liv
støver rundt i et arkiv
Læser i de gamle skrifter
selvom jeg får mange rifter
Jeg ved alt om Neros drifter
For historie er mit liv

\endverse
\beginverse
Jeg vil være pædagog
rende rundt i maosko
hjælpe børn i overtøjet
så de ikke bli'r tilmøjet
du kan tro de er fornøjet
jeg vil være pædagog

\endverse
\beginverse
Jeg vil være aktuar
hvor min skæbne den er klar
Forsikringsvidenskab er sagen
jeg vil ta' en bid af kagen
Mens jeg sylter kundeklagen
Jeg vil være aktuar

\endverse
\beginverse
Jeg vil muge ud hos køer
køre rundt med trillebør
sidde i en kæmpe traktor
moms det er en ukendt faktor
nu skal grisen ned til slagter
jeg vil muge ud hos køer

\endverse
\beginverse
Jeg vil læse på MI
for musik kan alle li'
Selvom jeg var sidst i klassen
Kan jeg stadig score kassen
blot jeg lær' at spille bassen
Men jeg kan ikke holde takten

\endverse
\beginverse
Jeg vil læse æstetik
i dansenes mystik
Svanens død den kan jeg lide
skønt jeg ej er en sylfide
vil jeg over scenen glide
Jeg vil læse æstetik

\endverse
\beginverse
Jeg vil læse nano tech
for min hjerne drak jeg væk
sidte fredag på Caféen?
da jeg rigtig sku' gå te'n
og jeg vælted om i sneen
jeg vil læse nano tech

\endverse
\beginverse
Jeg vil være datalog
og min fremtid den er god
gcc det er en gave
jeg blir nok en kodeslave
jeg vil gerne prøv' at snave
jeg vil være datalog

\endverse
\endsong



\beginsong{function fib(n)}[sr={Melodi:Sarah Brightman og Andrea Bocelli - Time to say goodbye}
,
by={DIKUrevy 2009}
,
cr={}]\hypertarget{function fib(n)}{}
\label{song10}

\beginverse
Læser man på DIKU skal man kunne programmere rekursioner
Man kan gøre det i Matlab kun med ganske få allokationer

\endverse
\beginverse
Let som en leg
Kod med mig, med mig

\endverse
\beginverse
Tænd for din laptop
Slå nu guien fra med \texttt{'matlab -nodesktop'}
Vi kan programmer'
og kalkuler'
Fibonaccis talrække

\endverse
\beginverse
\texttt{function fib(n)}
\hspace*{3mm}\texttt{p = 0;}
\hspace*{3mm}\texttt{A = 1}
\hspace*{3mm}\texttt{for j = 1 : n}
\hspace*{3mm}\hspace*{3mm}\texttt{temp = A;}
\hspace*{3mm}\hspace*{3mm}\texttt{A = A + p}
\hspace*{3mm}\hspace*{3mm}\texttt{p = temp;}
\hspace*{3mm}\texttt{end}

\endverse
\beginverse
og programmet er slut

\endverse
\endsong



\beginsong{Noget om helte}[sr={}
,
by={Tekst: Halfdan Rasmussen -- Melodi: Robert Normann}
,
cr={}]\hypertarget{Noget om helte}{}
\label{song11}

\beginverse
Livet er en morgengave,
sjælen er et pilgrimskor.
Der står krokus i min have,
der står øller på mit bord.
Under himlen hænger lærken
som et fjernt bevinget frø,
for en lærke tænker hverken
på at kæmpe eller dø.
\endverse

\beginverse
Her er fredeligt og stille,
her er ingen larm og støj.
Jeg har sået kruspersille,
og et brev med pure løg.
Lad alverden slå for panden
og bekæmpe spe med spot,
jeg vil enes med hinanden
og mig selv og ha' det godt.
\endverse

\beginverse
Samson gik og styrted' templer,
Peter Freuchen knak sit ben.
Ak mod disse to eksempler
er min dåd en sølle én.
Jeg har aldrig dræbt filistre
eller kæmpet med en haj,
og når stærke mænd blir bistre,
syn's jeg det er synd for mig.
\endverse

\beginverse
Der er nok som går og sysler
med at sprænge kloden væk.
Jeg vil ikke ha skærmydsler
og kanoner bag min hæk.
Mens de andre går og sveder
for at gi hinanden lak,
vil jeg pusle med rødbeder,
selleri og pastinak.
\endverse

\beginverse
Jeg vil ikke slås med bisser,
jeg vil så og ikke slå.
Selv de rødeste radiser
kan man roligt lide på.
Der er nok af danske helte
som er danske hele dag'n,
og går rundt og spænder bælte,
mens de råber; Fy for Fan.
\endverse

\beginverse
Tiden går og tiden hverver
store mænd til mandig dåd.
Jeg reserven blandt reserver
bryder ofte ud i gråd.
Jeg får nerver og migræne
blot jeg skær' mig på en dolk,
og vil hell're slå min plæne,
end slå løs på pæne folk.
\endverse

\beginverse
Livet er en dejlig gave,
jorden er en herlig jord.
der er øller i min mave,
der står krokus på mit bord.
Når reserverne skal stille
for at splitte kloden ad
skriver jeg med kruspersille
verdens mindste heltekvad.
\endverse
\endsong



\renewcommand{\thepage}{10}
\addtocounter{pageoffset}{-1}\beginsong{Fulbert og Beatrice}[by={},
,
by={}
,
cr={}]\hypertarget{Fulbert og Beatrice}{}
\label{song12}

\beginverse
I frankens rige, hvor floder rinde
som sølverstrømme i lune dal,
lå ridderborgen på bjergets tinde
med slanke tårne og gylden sal.
Og det var sommer med blomsterbrise
og suk af elskov i urtegård.
Og det var Fulbert og Beatrice,
og Beatrice var sytten år.

\endverse
\beginverse
De havde leget som børn på borgen,
mens Fulbert endnu var gangerpilt.
Men langvejs drog han en årle morgen,
mod Saracenen han higed' vildt.
Han spidded' tyrker som pattegrise,
et tusind stykker blev lagt på bår,
for Fulbert kæmped' for Beatrice,
og Beatrice var seksten år.

\endverse
\beginverse
Med gluttens farver på sølversaddel
han havde stridt ved Jerusalem.
Han kæmped' kækt uden frygt og dadel
og gik til fods hele vejen hjem.
Nu sad han atter på bænkens flise
og viste stolt sine heltesår,
som ganske henrykked' Beatrice
for Beatrice var femten år.

\endverse
\beginverse
En kappe prydet med små opaler
og smagfuldt ternet med tyrkens blod,
en ring af guld og et par sandaler
den ridder lagde for pigens fod.
Og da hun øjnede hans caprise
blev hjertet mygt i den væne mår.
Af lykke dånede Beatrice,
for Beatrice var fjorten år.

\endverse
\beginverse
Da banked' blodet i heltens tinding,
thi ingen helte er gjort af træ.
Til trods for plastre og knæforbinding
sank ridder Fulbert med stil på knæ.
Han kvad: ``Skønjomfru, oh skænk mig lise,
thi du alene mit hjerte rår!''
``Min helt, min ridder,'' kvad Beatrice,
for Beatrice var tretten år.

\endverse
\beginverse
Og der blev bryllup i højen sale
med guldpokaler og troubadour,
og under sange og djærven tale
blev Fulbert ført til sin jomfrus bur.
Og følget hvisked' om øm kurtice
og skæmtsom puslen blandt dun og vår.
For det var Fulbert og Beatrice,
og Beatrice var ell've år.

\endverse
\beginverse
Men ridder Fulbert den samme aften
af borgens sale blev båren død.
Den megen krig havde tær't på kraften,
og sejrens palmer den sidste brød.
Oh bejler, lær da af denne vise:
Ød ej din kraft under krigens kår.
Nej, spar potensen til Beatrice,
når Beatrice er født i går.

\endverse
\endsong

\newpage


\beginsong{Matematiksangen}[sr={Melodi: Svedbanken - Rollespilssangen}
,
by={Matematikrevyen 2014}
,
cr={}]\hypertarget{Matematiksangen}{}
\label{song13}

\beginverse
Det er ikke sjovt, så’n at træng’ til kaffe
og være løbet tør.
Hvor ku' jeg godt drik' en god kop kaffe,
men ingen kaffe til mig.
\endverse

\beginverse
Mine briller er grå og store,
og min ryg er krum.
Jeg har ik’ hørt om ham der Gustav,
så alle synes jeg er dum.
\endverse

\beginverse
Men når jeg’ så regner en uend'lig sum,
så bliver enhver der driller pluds'lig stum.
\endverse

\beginverse
For jeg læser mat’matik her på HCØ.
Jeg vil ikk’ holde op, så heller' dø.
Med ortonormale følger i Hilbertrum
får alle andre klø.
\endverse

\beginverse
Det er ikke sjov at folk syn’s jeg lugter
og siger min sang er grum.
Folk forstår ikk’ de sjove ordspil:
En aftenringnigs skum?
\endverse

\beginverse
Men når jeg’ så regner en uend'lig sum,
så bliver enhver der driller pluds'lig stum.
\endverse

\beginverse
Ja jeg læser mat’matik, mens jeg spiser kag’.
Du må kun være med, hvis du ka’ bag’.
Og jeg bruger theoremer, de viser vej.
 - så vis respekt for mig.
\endverse

\beginverse
Kom nu, bare en lille smule respekt.
Lyt nu - fordomme væk - se mig!
Kom nu!
\endverse

\beginverse
For jeg læser mat’matik her på HCØ.
Jeg vil ikk’ væk herfra, så hel’re dø.
Jeg læser mat’matik,
læs med mig, vær med, vær med, vær med.
\endverse

\beginverse
Jeg læser mat’matik!
Kom nu og skriv med kridt!! Skriv det! Skriv det! Kom med og skriv\ldots
\endverse
\endsong

\setcounter{temppage}{\value{page}}
\pagenumbering{shiftedpage}
\setcounter{page}{\value{temppage}}


\beginsong{Fuck mit liv}
[sr={Melodi: Bon Jovi -- It's My Life}
,
by={Matematikrevyen 2015}
,
cr={}]\hypertarget{Fuck mit liv}{}
\label{song14}

\beginverse
Dette er en emosang, yeah
Så pas nu på, for den kan godt bli' lang, yeah
Jeg har lige været ind' på KUnet Home
Ville melde mig til kurser, men så crashed' eduroam
\endverse

\beginverse
Fuck mit liv - og fuck mit studie
Ved ik' hvordan jeg ka' bru' det
Jeg vil meget heller' droppe ud
[Kor:] Drop nu uuuuud
Jeg trod' bar' vi sku' lægge sammen,
At mat'matik var fryd og gammen,
Dropper ud for byrden er massiv
Fuck mit liv!
\endverse

\beginverse
Jeg kan ikke finde ud af \TeX
Og studiegruppen - den er helt væk
I morgen bliver værre - der' lynopgav'
Jeg kan ikke printe - det skrev SCIENCE i et brev
\endverse

\beginverse
Fuck mit liv - og fuck mit studie
Ved ik' hvordan jeg ka' brug' det
Nu ska' jeg igen til reeksam'
[Kor:] Det' på Life
Har ikke fået forberedt mig
Glemt alle sider med notater
Frygter prøven bliver intensiv
Fuck mit liv!
\endverse

\beginverse
Lyt til mit råd - jeg er ikk' et udskud
Drop MATH, drop stress, baby, drop nu uuud!
\endverse

\beginverse
FUCK mit liv - og fuck mit studie
Ved ik' hvordan jeg ka' bru' det
Jeg vil bare hæve min SU
[Kor:] Du spilder klip
Kollegiet vil ha' mig opsagt
Er offer for de andres klapjagt (Kan ikke li' min emo-hårpragt)
Flytter hjem til mormors kollektiv
FUCK - MIT - LIV
\endverse

\beginverse
Og fuck mit studie
Jeg dropper ud - ka' ikke brug' det
Nu si'r jeg farvel til HCØ.
[Kor:] Kom ik' igen!
Ser frem til frihed og lidt opdrift
Men hov - hvad er nu det med fremdrift?
Tvangstilmeldt per rektors direktiv
FUCK - MIT - LIV!
\endverse
\endsong



\beginsong{Steve Hawking}[sr={Melodi: Postman Pat}
,
by={FysikRevy 2015}
,
cr={}]\hypertarget{Steve Hawking}{}
\label{song15}

\beginverse
Her er King
Steve Hawking
I sin stol, han er videnskabsidol

\endverse
\beginverse
Han dividér med nuller
Og skaber sorte huller
Hawking er en temmelig venlig mand

\endverse
\beginverse
Han fandt sin egen stråling
Det helt uden måling
Hawking er en meget særlig mand

\endverse
\beginverse
Han kører rundt hver dag
Som fysikgrøntsag
Han regner på sit fag
Og når han tænker måååske
Sker der vilde ting

\endverse
\beginverse
Det gjord' der ikk' i dag

\endverse
\beginverse
Her er King
Steve Hawking
Skrev en bog derfor er han skideklog

\endverse
\beginverse
Ved ikk' hvordan han gjorde
Men han kan stadig score
Hawking må ha' virk'lig meget SWAG

\endverse
\beginverse
Hawking han har helt vildt meget SWAG

\endverse
\beginverse
Hawking han har helt vildt meget SWAG

\endverse
\endsong



\beginsong{Thinkpad T40}[sr={Melodi: Roger and Over - Volvo B18-210}
,
by={DIKUrevy 2014}
,
cr={}]\hypertarget{Thinkpad T40}{}
\label{song16}

\beginverse
Vi så hinanden første gang
i Store UP(1)
Du var slukket, jeg var fuld -
din ejer han var væk
Gik en lille tur med dig,
i Kantinen var der plads
Jeg satte mig, og tændte dig
fik stød og det gjord' nas

\endverse
\beginverse
Du er min Thinkpad T40, og du er slidt
Du mangler lidt RAM, og du lugter lidt af kridt
68 graders kerne, hver gang du starter op
Og dit lunefulde drev, før mig til at ta' backup

\endverse
\beginverse
Du var en gammeldags maskine
Det med 'nyt' det var lidt svært
Men du holdte mig allig'vel
Underholdt med lidt af hvert
Du havde kræfter til at spille
gamle spil fra MS-DOS,
de kalder dig en bærbar
men du er sgu lidt en klods

\endverse
\beginverse
Du er min Thinkpad T40, og du er slidt
Du mangler lidt RAM, og du lugter lidt af kridt
68 graders kerne, hver gang du starter op
Og dit lunefulde drev, får mig til at ta' backup

\endverse
\beginverse
Du er min Thinkpad T40, og du er slidt
Du mangler lidt RAM, og du lugter lidt af kridt
68 graders kerne, hver gang du starter op
Og dit lunefulde drev, får mig til at ta' backup

\endverse
\beginverse
Vi skrev bachelor sammen,
med gcc og vim
Når du laved' sjov, og sletted' alt,
så skrev jeg det igen
Det var dagen før min deadline,
jeg sku' oversæt' min TeX
Jeg låste dig, men ik' min dør
og pludselig var du væk

\endverse
\beginverse
Du var min Thinkpad T40, og du var slidt
Du mangled' lidt RAM, og du lugted' lidt af kridt
68 graders kerne, hver gang du started' op
Og dit lunefulde drev, fik mig til at ta' backup

\endverse
\endsong



\beginsong{En ringe historie}[sr={Melodi: Se den lille kattekilling}
,
by={Matematikrevyen 2012}
,
cr={}]\hypertarget{En ringe historie}{}
\label{song17}

\beginverse
Betragt en lille abelsk ring -- bemærk, den kommuter'.
Den har netop ét 0 og ét 1 og ikke mer'.
Se, nu går den hen til $\mathbb C$ og indleder med ``Hej''.
``Skal vi være isomorfe bare du og jeg?''
``Nej, det vil jeg ikke. Du er tællelig min ven
-- ikke på den seje måde som min homie $\mathbb N$.''
Ingen ville lege med en ring der' triviel,
så den måtte indgå automorfisk med sig selv.
\endverse
\endsong



\beginsong{Mat'matik}[sr={Melodi: tema fra Bubbi-bjørnene}
,
by={FysikRevy 2008}
,
cr={}]\hypertarget{Mat'matik}{}
\label{song18}

\beginverse
Vektorfunktioner i tre dimensioner
skal opereres med curl og divergens.
$\nabla^2$ gir koldsved på panden;
planintegraler er en pestilens.
\endverse

\beginverse
Mat'matik
strider ofte mod enhver logik;
men har du elektrodynamik,
så skal du ku' mat'matik.
\endverse

\beginverse
Kan en hermitisk og injektiv matrix
ha' negativ trace men en nul-determinant?
Solovej siger dens egenværdier
vil være reelle, men er det mon sandt?
\endverse

\beginverse
Mat'matik
strider ofte mod enhver logik;
men hvis du har kvantemekanik,
så skal du ku' mat'matik.
\endverse

\beginverse
$T_a$ vil gi' $SU(3)$-symmetri,
og Lagrange-operator'n er Gauge-invariant.
Når operator'ne virker på kvarkerne,
får de en farve og smagen af kvant
\endverse

\beginverse
Mat'matik
strider ofte mod enhver logik;
men i kvantekromodynamik,
så skal du ku' mat'matik.
\endverse
\endsong



\beginsong{Papirsklip}[sr={}
,
by={Kim Larsen}
,
cr={}]\hypertarget{Papirsklip}{}
\label{song19}

\beginverse
Når nu min verden bliver kold og forladt
finder jeg trøst i min kæreste skat
klipper motiver med drømme og saks
papirsilhuetter af den fineste slags.
\endverse

\beginverse
Her er et af min far og min mor
de som gav mig til denne jord
kærlige kys og en duft af jasmin
altid solskin og sødeste min.
Livet er langt
lykken er kort
salig er den der tør give det bort.
\endverse

\beginverse
Her er et af min bedste ven
utallige klip går igen og igen
et aldrig færdigt eller fuldendt motiv
sort silhuet af kvinden i mit liv.
Livet er langt\ldots
\endverse

\beginverse
Domine et sanctus
Domine et sanctus
Domine et sanctus
\endverse

\beginverse
Når nu min verden bliver kold og forladt
finder jeg trøst i min kæreste skat
klipper motiver med drømme og saks
papirsilhuetter af den fineste slags.
\lrep Livet er langt
lykken er kort
salig er den der tør give det bort.\rrep
\endverse
\endsong



\beginsong{Jobsøgning}[sr={Melodi:Klokkeren fra Notre Dame - Hell fire}
,
by={DIKUrevy 2013}
,
cr={}]\hypertarget{Jobsøgning}{}
\label{song20}

\beginverse
Hør her, kære venner,
I ved at jeg er helt speciel,
I ved at jeg aldrig koder slam

\endverse
\beginverse
Hør her, kære venner,
I ved at jeg er bedre end
En gennemsnitlig Java-programmør

\endverse
\beginverse
Han imponerer!

\endverse
\beginverse
Så sig mig nu, venner,
Hvorfor har jeg intet job,
Når underkoderne bli'r revet væk?

\endverse
\beginverse
Vi ved det ikke!

\endverse
\beginverse
Hvornår vil de lære
At fremtiden er funktionel,
At PHP og Python snart vil dø!?

\endverse
\beginverse
Så kom og kod for ooos!

\endverse
\beginverse
Med Haskell
Og Erlang
Bliver jeg overset
De job jeg
Begærer
Bli'r aldrig mig til del...

\endverse
\beginverse
Det' ik' min skyld

\endverse
\beginverse
Det ik' min fejl!

\endverse
\beginverse
Vi har kaffe!

\endverse
\beginverse
Det' kapitalens valg af sprog
der gør det svært!

\endverse
\beginverse
  frem.

\endverse
\beginverse
Du kan også få pizza!

\endverse
\beginverse
Det' ikk min skyld,

\endverse
\beginverse
Hvad med bacon?

\endverse
\beginverse
Hvis kapital'n

\endverse
\beginverse
Eller kage?

\endverse
\beginverse
Går konkurs fordi de sagde nej til mig!

\endverse
\beginverse
Sig nu ja, vi vil ha' dig!

\endverse
\beginverse
Hvorfor skal man bare
Acceptere dårlig stil,
Finde sig i BlueJ og Eclipse?

\endverse
\beginverse
Og ofre sin sjæl på
Javas alter for et job?
Jeg vil aldrig kode OOP!

\endverse
\beginverse
Øh, Rasmus? Google har ringet - IGEN.

\endverse
\beginverse
HVAD?!

\endverse
\beginverse
De mangler altså virkelig en Haskellprogrammør...

\endverse
\beginverse
For pokker... Så sig at jeg kun koder... øh... ALGOL! Og smut så.

\endverse
\beginverse
Med ALGOL
Og Agda
Holder jeg mig obskur.
De fanger
Mig aldrig
Undslipper gang på gang!

\endverse
\beginverse
Vi giver snart helt op.

\endverse
\beginverse
Hvis de vidste besked...

\endverse
\beginverse
Vil du have det job?

\endverse
\beginverse
Hvis de bare indså...

\endverse
\beginverse
Han får aldrig et job.

\endverse
\beginverse
Jeg kan ikke kode
NOGET
SOM
HEEEELST!

\endverse
\endsong



\beginsong{$\lambda$-kalkylen}%21
[sr={Melodi: The Llama Song}
,
by={DIKUrevy 2011}
,
cr={}]\hypertarget{$lambda$-kalkylen}{}
\label{song21}

\beginverse
Her et lambda, dér et lambda
Og et ganske lille lambda
ML-lambda, Haskell-lambda
Lambda, lambda
Fail!

\endverse
\beginverse
Æg og alligator-lambda
Java har slet ikke lambda
Lambda i .NET, Omega
Lambda, lambda
Fail!

\endverse
\beginverse
Jeg har fået value
og drukket meget te
Jeg har lært om Hitler
Dinosaurer og VB
Jeg har kodet uafbrudt i 42 døgn
Men så døde jeg af sult
for kagen var en løgn

\endverse
\beginverse
Har du nog'nsinde set et lamdba
i et lambda, på et lambda
Lambda x dot færdig med x
Lambda, lambda
Fail!

\endverse
\beginverse
Om nom nom nom nom nom lambda
Og et A forklædt som lambda
Lambda i \LaTeX, Omega
Lambda, lambda
Fail!

\endverse
\beginverse
Jeg har kodet Brainfuck
og dyrket lidt motion
Jeg har været udenfor
det var en perversion
Jeg har skrevet kode fyldt med map og head og tail
Nu vil jeg blive kandidat
og gå ned med fail.

\endverse
\endsong



\beginsong{Forever DIKU}[sr={Melodi: Rick Astly - Never Gonna Give You Up}
,
by={DIKUrevy 2015}
,
cr={}]\hypertarget{Forever DIKU}{}
\label{song22}

\beginverse
	    Jeg er startet i år
            Du' ikke en rus, men det er jeg
            Jeg tager kurser, håber jeg består
            Til mange andre, der -- giver du et nej

\endverse
\beginverse
            Om mig ser jeg folk der dropper studiet
            (Men) jeg vil ej for-lade dig

\endverse
\beginverse
            Jeg vil aldrig droppe ud
            Jeg vil aldrig stå for skud
            Jeg vil aldrig få et job og bar' stoppe
            Jeg vil aldrig sige op
            Jeg vil aldrig følge trop
            Jeg vil aldrig helt få nok af DIKU

\endverse
\beginverse
            % Jobtilbud
            Tiden den løber, jeg får job
            En gang om ugen men
            De vil ha' mig mere
            Hvis det er sådan må jeg -- sige op
            De andre dage skal -- jeg studere

\endverse
\beginverse
            Og hvis jeg så får brug for penge
            Ta'r jeg et in-struktorat

\endverse
\beginverse
            Jeg vil aldrig droppe ud
            Jeg vil aldrig stå for skud
            Jeg vil aldrig få et job og bar' stoppe
            Jeg vil aldrig sige op
            Jeg vil aldrig følge trop
            Jeg vil aldrig helt få nok af DIKU

\endverse
\beginverse
            Jeg vil aldrig droppe ud
            Jeg vil aldrig stå for skud
            Jeg vil aldrig få et job og bar' stoppe
            Jeg vil aldrig sige op
            Jeg vil aldrig følge trop
            Jeg vil aldrig helt få nok af DIKU

\endverse
\beginverse
            (Ååh, drop nu ud)
            (Ååh, drop nu ud)
            Aldrig droppe ud, aldrig droppe ud
            (Drop nu ud)
            Aldrig droppe ud, aldrig droppe ud
            (Drop nu ud)

\endverse
\beginverse
            % Graduation
            Men studietiden er så kort
            Har været flittig så,
            mangler kun specialet
            Jeg kan ej bære at sku' drage bort
            Efter mit forsvar, så -- hej for evigt?

\endverse
\beginverse
            Kan jeg ikke få en ekstra omgang?
            Jeg må tag' en Ph.D.

\endverse
\beginverse
            Jeg vil aldrig droppe ud
            Jeg vil aldrig stå for skud
            Jeg vil aldrig få et job og bar' stoppe
            Jeg vil aldrig sige op
            Jeg vil aldrig følge trop
            Jeg vil aldrig helt få nok af DIKU

\endverse
\beginverse
            Jeg vil aldrig droppe ud
            Jeg vil aldrig stå for skud
            Jeg vil aldrig få et job og bar' stoppe
            Jeg vil aldrig sige op
            Jeg vil aldrig følge trop
            Jeg vil aldrig helt få nok af DIKU

\endverse
\beginverse
            Jeg vil aldrig droppe ud
            Jeg vil aldrig stå for skud
            Jeg vil aldrig få et job og bar' stoppe
            Jeg vil aldrig sige op
            Jeg vil aldrig følge trop
            Jeg vil aldrig helt få nok af DIKU

\endverse
\endsong



\beginsong{Faster than Light}[sr={Melodi:Backstreet Boys - Larger than Life}
,
by={FysikRevy  2012}
,
cr={}]\hypertarget{Faster than Light}{}
\label{song23}

\beginverse
We are working hard at il OPERA experiment
A pretty student comes and she looks so magnificent
blonde hair, stilleto, and soft skin so bright

\endverse
\beginverse
Small neutrinos - hard to see, hard to see
we need to determine their velocity
(to) concentrate is hard (with) pretty girls in sight
And that makes them faster than light

\endverse
\beginverse
Gorgeous bella donna make us drop our activity
Now we have refuted general relativity
blonde hair, stilleto, soft skin so bright

\endverse
\beginverse
Small neutrinos - hard to see, hard to see ...

\endverse
\beginverse
blonde hair, stilleto, soft skin so bright

\endverse
\beginverse
Small neutrinos - hard to see, hard to see ...

\endverse
\endsong



\beginsong{Her på mat'matik}
[sr={Melodi: Fætter Mikkel}
,
by={Matematikrevyen 2015}
,
cr={}]\hypertarget{Her på mat'matik}{}
\label{song24}

\beginverse
Her på mat'matik
Skal man være kvik
Sætte pris på elegance
Ha' intuition
Også ambition
Gribe fat i hver en chance
\endverse

\beginverse
For vi elsker ren logik (klap-klap)
algebra og statistik (klap-klap)
Alt der er komplekst
Analyse, vækst
det er her vi har det bedst (klap-klap)
\endverse

\beginverse
Fra det første år
Kurserne består
Fællesskabet ej forglemme
Der bli'r undervist
Masser kage spist
Det er her hvor vi har hjemme
\endverse

\beginverse
For vi elsker ren logik (klap-klap)
algebra og statistik (klap-klap)
Alt der er komplekst
Analyse, vækst
det er her vi har det bedst (klap-klap)
\endverse


\endsong



\beginsong{Censor Baby}[sr={Melodi: Eartha Kitt - Santa Baby}
,
by={MBKrevy 2014}
,
cr={}]\hypertarget{Censor Baby}{}
\label{song25}

\beginverse
Censor baby, så er vi til eksamen igen, min ven
Jeg har arbejdet hårdt
Censor baby, så giv mig det tolv-tal i dag

\endverse
\beginverse
Censor baby, og hvilket spørgsmål holder du af, i dag?
Lad mig svare på det
Censor baby, så giv mig nu det tolv-tal i dag

\endverse
\beginverse
Tænk nu på min karakter
Se hvordan jeg virk'lig gerne vil ha' mer'
Jeg vil helst ikke ta' den om
Men under 12 er en unfair dom

\endverse
\beginverse
Censor baby, nu kender jeg min spektroskopi, fordi
Jeg har studeret det hårdt
Censor baby, så giv mig nu det tolv-tal i dag

\endverse
\beginverse
Censor sveske, jeg synes, at du er et geni, ka' du li'
En molbos anatomi
Censor cutie, så giv mig nu det tolv-tal i dag

\endverse
\beginverse
Spørg om alt, jeg skulle lære
Men skal vi ikke sige, at du lader være
Består du mig, så vil jeg gi
Privatlektion i blomst og bi

\endverse
\beginverse
Flænser baby, jeg tror censor vil sige et ord
Jeg tror, at du skal lytte til ham
Flænser baby, votering er en to-sidet sag

\endverse
\beginverse
Se hvor har jeg været sød
Tænk på al' de Fernet shots jeg ikke nød
Når du streger med din pen
Så husk, jeg er din bedste ven

\endverse
\beginverse
Censor baby, jeg mangler kun at nævne en ting, Latex-s
Ikke kun kodning for mig

\endverse
\beginverse
Censor baby, så giv mig nu det tolv-tal i dag
Giv mig nu det tolv-tal i dag
Et tolv-tal...i nat

\endverse
\endsong



\beginsong{Fag Facisme}[sr={Melodi:Julia Ward Howe - The Battle Hymn of the Republic}%26
,
by={FysikRevy  2013}
,
cr={}]\hypertarget{Fag Facisme}{}
\label{song26}

\beginverse
Så tak os for at masse faktisk oss' er energi
og tak os for konstant forøgelse af entropi
og tak os for at tyngdekraften virker overalt
for vi følger Newtons kald

\endverse
\beginverse
Mat'matik approksimeres
Datalogen ignoreres
Humanister likvideres
Ren Fysik det bli'r vor svar!

\endverse
\beginverse
Så tak os for at verd'nen er en superposition
og tak os for at elektronen er en fermion
og tak os for at kvantemekanik er mer' end ord
Udtænkt af vor Niels Bohr

\endverse
\beginverse
Mat'matik approksimeres ...

\endverse
\beginverse
Vi diff'rentierer step-funktioner uden megt' besvær
Vi indekserer kun med 1 og nester for-loops her
På KUA har I kommunister, dræn og meta-vrøvl
Kom an, vi gir' jer høvl

\endverse
\beginverse
Mat'matik approksimeres ...

\endverse
\beginverse
Så tak os for at alting overhoved' eksister',
og tak os for de svære ting vi dagligt integrer'.
Og tak os for at sætte flotte streger gennem $h$
ren fysik det vil bestå

\endverse
\beginverse
Mat'matik approksimeres ...

\endverse
\beginverse
Mat'matik approksimeres ...

\endverse
\endsong



\beginsong{Ode til Analyse 0}%27
[sr={Melodi: ABBA -- Gimme! Gimme! Gimme! (Man After Midnight)}
,
by={Matematikrevyen 2015}
,
cr={}]\hypertarget{Ode til Analyse 0}{}
\label{song27}

\beginverse
Når du skal
Op til An0 eksamen, er det ikke forkert
At ha' styr på uniform kontinuert
Hvis du står
Og skal til at forklare lidt om hovedsætning 3
At bevise den det er jo ingen sag
\endverse

\beginverse
Når slaget går igang
Så husk blot denne sang
\endverse

\beginverse
\emph{Når vi, når vi, når vi har epsilon givet
Så vær' ikke bange, for der' noget du ved
Når vi, når vi, når vi har epsilon givet
Findes der et delta at parere det med}
\endverse

\beginverse
Skal du til
At bevise Bolzano, eller Riemanns integral
(Vis) Begge to på samme tid helt uden pral
Husk da bare
At så får du brug for middelsummer og et I
Trekantsulighed det' noget vi ka' li'
\endverse

\beginverse
Når slaget går igang
Så husk blot denne sang
\endverse

\beginverse
\emph{Når vi, når vi, når vi\ldots}\rep{2}
\endverse

\beginverse
Når slaget går igang
Så husk blot denne sang
\endverse

\beginverse
\emph{Når vi, når vi, når vi\ldots}\rep{2}
\endverse


\endsong



\beginsong{Ode til Analyse 1}%28
[sr={Melodi: Shout -- Jeg tror det kaldes kærlighed}
,
by={Matematikrevyen 2013}
,
cr={}]\hypertarget{Ode til Analyse 1}{}
\label{song28}

\beginverse
Hvordan gi’r vi mening til
Et uend’ligt summespil?
Tallene bli’r fler og fler
Tjek om delsum konverger’

\endverse
\beginverse
Sum af $n$’te del i $p$
Forholdstest og rottesten
Vi lærte om en dejlig ven
Den hedder Grænsesammenligningstesten

\endverse
\beginverse
\emph{[Omkvæd:]
Denne sang den er til dig
Åh Jan Philip Solovej
Når funktionerne de ændrer sig
Så ta’r vi det skridt for skridt
Lært i Analyse 1
Lad $n$ gå mod uendeligt}

\endverse
\beginverse
Nu blev faget mer’ abstrakt
Vi fandt en sum for $\pi$ eksakt
Lærte om vor tavlesvamp
Valget tør og våd er\ldots ikke mer’ en kamp

\endverse
\beginverse
\emph{[Omkvæd]}

\endverse
\beginverse
Er du nu lang langt væk eller så nær?
(Er du nu lang langt væk eller så nær?)

\endverse
\beginverse
Er metrikken 0, så du er her?

\endverse
\beginverse
Wowyeahhh

\endverse
\beginverse
\emph{[Omkvæd]}\rep{2}

\endverse
\beginverse
\emph{Lad $n$ gå mod uendeligt}

\endverse
\endsong



\beginsong{Hvad har du lavet på din datamat?}[sr={Melodi: Cirkusrevyen 1967 - Hvem har du kysset i din gadedør?}
,
by={DIKUrevy 2015}
,
cr={}]\hypertarget{Hvad har du lavet på din datamat?}{}
\label{song29}

\beginverse
Jeg er 30 år og bachelor,
her på data( - )logi
Vores log har set du brugte Tor,
her til morgen, klokken ni.
Hr. Betjent, det må man gerne
Har du måske glemt din hjerne?

Hvad har du (jeg) lavet på din (min) datamat,
i dit (mit) hackeri?
Jeg har blot øvet på min kandidat
Det kan jeg godt li'!

Næh, du vil bare surfe ud'n kontrol
Så du kan se på død og vold
Fodre et kæledyr med alkohol
Og andet svineri!

Kan du så gå din vej med dit forhør?
SÅ unødvendig'

\endverse
\beginverse
Hvad har du lavet på din (min) datamat
I dit hackeri?!
Fyyyyh..!
Nej, men du forstår slet ikke hvad jeg har gang i.
Det er en BOMBE under den nuværende infrastruktur.
jeg kommer til at sprænge muren for effektivitet!

\endverse
\beginverse
Du skal ikke bombe noget som helst
Det er regulær terrór
Hr. Betjent, du misforstår mig vidst
Det var blot en metafor!
Nej, men det kan vær' det samme
Metaforer de er klamme!

\endverse
\beginverse
Hvad har du lavet på din datamat
I dit hackeri?
Hvem så dig sidste ug' i Labitat?
Det gjord' politi!
Det er måske en... naaaarko-pizza?
  Jeg tror at du har hentet programmel
  fra www.kriminel
  Den slags man bruger til at slå ihjel
  Oh fy, oh føj, forbudt!

\endverse
\beginverse
  Deeeeeeeeet har man bestemt derned' i Bruxelles
  Fine lille snut!
  Hvad har du (jeg) kodet på din datamat?
  I dit (mit) hackeri?

\endverse
\beginverse
Det er måske en... naaaarko-pizza?
  Jeg tror at du har hentet programmel
  fra www.kriminel
  Den slags man bruger til at slå ihjel
  Oh fy, oh føj, forbudt!

\endverse
\beginverse
  Deeeeeeeeet har man bestemt derned' i Bruxelles
  Fine lille snut!
  Hvad har du (jeg) kodet på din datamat?
  I dit (mit) hackeri?

\endverse
\beginverse
Du er sand'lig haft det ganske nemt,
Jeg har fikset Rejsekort
Du har forbrudt dig på danske softwarepatenter. Det må man ikke!
IH! DET.. DET ER IHVERTFALD BARE MEGA ULOVLIGT!
Hr. Betjent, du må vær' ra-ar
Lad mig bare kort forklare.
Det, som jeg laved' på min datamat,
i mit hackeri
var blot at fixe hele Rejsekort,
så, tag du bare fri

\endverse
\beginverse
Men den slags rebeller kan vi ikke li'
I vores demo-bureaukrati,
og inden at du har fået talt til ti
Da bli'r du sat fast!
Ingen' dataloger de må røre ved
Statens kort af plast!

\endverse
\beginverse
  Så nu er du (jeg) anholdt i dit (mit) hackeri
  Du' under arrest!

\endverse
\beginverse

\endverse
\endsong



\beginsong{Integralsangen}%30
[sr={Melodi: Peberkagebagesangen}
,
by={Matematikrevyen 2008}
,
cr={}]\hypertarget{Integralsangen}{}
\label{song30}

\beginverse
Hvis\ldots
integralet du skal løse,
må du ikke bare sløse.
Du skal give det dit bedste,
og så ikke blive gal.
Og hvis stamfunktionen findes,
er det bare om at mindes:
Differens af indsat grænser
det er vores areal.

\endverse
\beginverse
Hvis du ikke kender grænser,
differens du ikke ænser;
du skal find’ de tal der matcher
de værdier som du har.
Hvis konstanten ellers stemmer,
vores resultat vi kender.
Din instruktor bliver glad,
og jeres opgave er klar.

\endverse
\endsong



\beginsong{Når en humanist adderer}%31
[sr={Melodi: Peberkagebagesangen}
,
by={TÅGEKAMMERET}
,
cr={}]\hypertarget{Når en humanist adderer}{}
\label{song31}

\beginverse
Når en humanist adderer,
tar’ hun først et $\pi$ i fjerde,
ganger det med $n$ matricer,
mens hun vælger en ``sød'' brøk.
Læg det til determinanten
det er humanistisk fjanten -–
for hun når jo aldrig læng’re
end at svaret det er $x$.
\endverse

\beginverse
Når en biolog skal lære,
hvordan man multiplicerer,
skal eksempler observeres,
hvor to køer bli’r til tre.
Og gør man det en gang mere,
Så bli’r tre kø’r jo til flere!
Så hun kender kun’ et svar:
At der er flere nu end før.
\endverse

\beginverse
Geologer de kan regne,
dem skal man ikke forklejne.
De bestemmer snildt en alder,
når de måler på en sten.
Men har deres kompetencer
nogen sociale nuancer?
Man skal være i en grusgrav
før det bliver relevant!
\endverse

\beginverse
Men hos os fra TÅGEKAMMERET
giver regning ingen jamren,
vores ligninger gi’r mening
også med socialt aspekt.
Vi beregner øl i kasser
med potens og inderklasser,
og vor KA\$\$ har altid styr på,
hvem der skylder hvad for hvad.
\endverse

\beginverse
Så lektionen den må være:
Hold mat/fys’erne i ære,
de har styr på deres sager,
og de drikker mange øl. [Skål!]
Udregningerne skal komme
fra han-dyrene med vomme,
der beviser de har været
alt for meget på TK!
\endverse
\endsong



\beginsong{Linieskriverdriversangen}%32
[sr={Melodi:Dyrene i Hakkebakkeskoven "Peberkagebagesangen"}
,
by={DIKUrevy 2005}
,
cr={}]\hypertarget{Linieskriverdriversangen}{}
\label{song32}

\beginverse
Når en skriverdriverskriver
skriver linieskriverdriver
skal en skriverdriver skrive
linieskriverdrivprogram.

\endverse
\beginverse
For det hedder ikke driver
nej, det' bare nog't man skriver
men bar' fordi Microsoft de skriver
driver bliver driver ikke nødvendigvis korrekt, vel?

\endverse
\beginverse
Hvis en lagerlærer lærer
andre lagerlærer's lære
bør der's lagerlærdom lagres
i et lærelagerlag.

\endverse
\beginverse
Men det er jo så'n med lærer
at de ik' ka' lagre lære
for de kager rundt og tager
jo slet ikke hinanden alvorligt, vel?

\endverse
\beginverse
Når processor'n processerer
og processerne bli'r flere
og professor'n ik' ka' 'pere
fler' processer i hans RAM.

\endverse
\beginverse
{colb}
Ja, så må vi allokere
fler' studerende med mere
lager end der's lagerlærer
har i lagerlærerlag.

\endverse
\beginverse
Ja, så må vi allokere
fler' studerende med mere
tid til at studere negre
så de ka' forny mit professorat! Ikke sandt?

\endverse
\beginverse
Når din vektorlektor nægter
at en sektorvektorhægter
peger på en sektorvektor
er din vektorlektor gal.

\endverse
\beginverse
For en sektorvektorhægter
hægter kun til sektorvekt're
og det' nytter ik' at lektor
nægter hægters eksistens.

\endverse
\beginverse
For en sektorvektorhægter
hægter kun til sektorvekt're
selvom alle vor's indtægter
mistes når en vektorlektor møder vektorlektorinspektoren, ikke?

\endverse
\beginverse
Hvis din pegerpeger peger
På noget andet end en peger
Så' din pegerpegers pegning
Gået hen og blevet gal

\endverse
\beginverse
For en pegerpeger plejer
Kun at peg' på andre peger
Så vi retter pegerpeger'n
For at den kan bli' normal

\endverse
\endsong



\beginsong{Tål daj i læseferien}[sr={Melodi: 12 days of Christmas}
,
by={Mona Holm og Marianne Bangsø, Julerevy 1987}
,
cr={}]\hypertarget{Tål daj i læseferien}{}
\label{song33}

\beginverse
Husk: der skal skåles efter 4. vers -- hver gang!
\endverse

\beginverse
På den 1. da’ i læseferien sa’ jeg til mig selv:
Der er god tid, jeg pjækker i dag.
\endverse

\beginverse
På den 2. da’ i læseferien\ldots
Mat’matik er trist.
\endverse

\beginverse
På den 3. da’ i læseferien\ldots
Jeg er træt.
\endverse

\beginverse
På den 4. da’ i læseferien\ldots
Tror jeg går på druk.
\endverse

\beginverse
På den 5. da’ i læseferien\ldots
Åh, tømmermænd.
\endverse

\beginverse
På den 6. da’ i læseferien\ldots
Der er fest i aften.
\endverse

\beginverse
På den 7. da’ i læseferien\ldots
Hvor er min bog nu.
\endverse

\beginverse
På den 8. da’ i læseferien\ldots
Der' film i TV.
\endverse

\beginverse
På den 9. da’ i læseferien\ldots
Savner mine venner.
\endverse

\beginverse
På den 10. da’ i læseferien\ldots
Mon jeg kan nå det.
\endverse

\beginverse
På den 11. da’ i læseferien\ldots
Nu skal der læses.
\endverse

\beginverse
På den 12. da’ i læseferien\ldots
Gud, det' i morgen.
\hspace{0em}\phantom{M}\vdots
Mat’matik er trist.
Så'n gik tiden -- jeg dumped' i dag!
\endverse
\endsong



\beginsong{Bust'ed}[sr={Melodi: Buster}
,
by={Erik Søe Sørensen og Troels Bjerre Sørensen, Majrevy 2003}
,
cr={}]\hypertarget{Bust'ed}{}
\label{song34}

\beginverse
Stille, stille, stille, jeg har tømmermænd
og min hjerne er vist svundet
det sidste jeg så var en lunken Top
som blev for hurtigt bundet.

\endverse
\beginverse
Det hyler og tuder i min hjerne nu
væk er spiritussens charmer
på gaden kører bilerne hurtigt forbi
og fuglene de larmer.

\endverse
\beginverse
Jeg’r busted
åh, busted
Jeg har brækket mig i cisternerne
og ønsker mig langt ud i det blå.

\endverse
\beginverse
Stille, stille, stille, jeg har tømmermænd
der er larm, selv under dynen
jeg strækker mig efter en doven Top
og putter den i trynen

\endverse
\beginverse
I aftes gjord’ jeg noget som man ikke skal
man skal nemlig ikke tylle
til sidst sag’ jeg dav til en stor og mægtig spand
som jeg næsten ku’ fylde

\endverse
\beginverse
Jeg’r busted
åh, busted
Jeg tog vist en pige på lanternerne
min kind er nu langt ud i det blå

\endverse
\beginverse
Jeg’r busted
åh, busted
Som sprunget ud af et mareridt
tager tømmeren hævn, og han kræver sit
åh-nej\ldots
Jeg ønsker ham langt ud i det blå
Jeg’r busted – åh, busted.

\endverse
\endsong



\beginsong{Facebookevents}[sr={Melodi:TV2 - Det Mørke Jylland}
,
by={DIKUrevy 2014}
,
cr={}]\hypertarget{Facebookevents}{}
\label{song35}

\beginverse
Jeg tjekker mails, der er intet nyt
Jeg er ikk' inviteret
For jeg er nemlig, slet ikk' på facebook
Men jeg vil stadig med

\endverse
\beginverse
Før i tiden, der fik man en mail
Når der var fest og sjov
Og uden facebook, der ved jeg ikk'
Hvornår der er gilde

\endverse
\beginverse
Jeg ved at jeg sandsynligvis går glip af noget
Og tænk hvis jeg nu ikke får det hele med

\endverse
\beginverse
På facebook hav-de jeg tilmeldt mig
et så vildt event
Det var fedt, og jeg havd' glædet mig
Men jeg kom ikk' af sted

\endverse
\beginverse
(For) jeg sku' lige ha' tjekket min væg
Og like hvad jeg så
Havde få't notifikation
Sku' lige poke en ven

\endverse
\beginverse
Jeg ved at jeg sandsynligvis går glip af noget
Og tænk hvis jeg nu ikke får det hele med
Der gik livet lig' forbi med klar besked
Mens jeg selv var væk et øjeblik
Et helt andet sted

\endverse
\beginverse
Jeg sidder her til mit eget party
Men hvor bli'r gæsterne af?
De meldte sig alle til mit event
Men logged' ah-aldrig af

\endverse
\beginverse
Der gik livet lig' forbi med klar besked
Mens jeg selv var væk et øjeblik
Et helt andet sted
Der gik livet lig' forbi med klar besked
Mens jeg selv var væk et øjeblik
Et helt andet sted

\endverse
\endsong



\beginsong{diff}[sr={Melodi:Aqua - Doctor Jones}
,
by={DIKUrevy 2013}
,
cr={}]\hypertarget{diff}{}
\label{song36}

\beginverse
Ofte er filerne ens
Lige nu er de det ikke
Hvem har mon rørt min C-fil?
Hvad har de indsat i koden?

\endverse
\beginverse
\texttt{void public static}
\texttt{void main}, wat?
\texttt{void abstract class}, øv!

\endverse
\beginverse
Men min kode er stor,
den kan jeg ikk' oversku'
Mangler den rigtige mand,
en som kan redde mig nu

\endverse
\beginverse
R: Hey baby, hva' mangler du
S: At finde forskelle
R: Ska' det være lige nu?
S: Kom nu og \texttt{diff}
R: Jeg kigger på din' filer nu
S: Hva' sker der når du ser?
R: Jeg kan nu ingen forskel se, pyt, det kan jo ske

\endverse
\beginverse
René Diff, Diff,
Tjekker filerne.
René Diff, René Diff,
Er de ens!
De er ens!

\endverse
\beginverse
René Diff, Diff,
filernes sherif.
René Diff, René Diff,
Er de ens!
De er ens!

\endverse
\beginverse
\texttt{void public static} ...

\endverse
\beginverse
\texttt{void public static} ...

\endverse
\beginverse
Er der klasser i C?
Nej, det er da vist i Java
René, kan du ikk' se?
Du er et levn fra halvfemserne

\endverse
\beginverse
R: Hey baby, hva' er det nu
S: Dig kan jeg ikke brug'
R: Lad mig tjekke den her, du
S: Det er en \texttt{.gif}!
R: (host) det er et billed' af en grif
S: Du var min diff-kalif!
R: Hvad vil' du gøre uden mig, jeg er da for sej

\endverse
\beginverse
René Diff, Diff,
lad mig gi' et fif
René Diff, René Diff,
Styrtflyver!
Du lyver!

\endverse
\beginverse
René Diff, Diff,
synger næste riff
René Diff, René Diff,
Hvad vil du!
Ha' ret nu!

\endverse
\beginverse
\texttt{void public static} ...

\endverse
\beginverse
\texttt{void public static} ...

\endverse
\beginverse
Brug'r nu diff3
Brug'r nu diff3
Brug'r nu diff3
Brug'r nu diff3

\endverse
\beginverse
René Diff, Diff,
vågn op nu
René Diff, Diff,
forstå nu
René Diff, Diff,
du er sgu
René Diff, Diff,
i udu

\endverse
\beginverse
\texttt{void public static} ...

\endverse
\beginverse
\texttt{void public static} ...

\endverse
\beginverse
René Diff, Diff
Giv min fil et snif
René Diff, René Diff
Han tog fejl
Det gjor' jeg

\endverse
\beginverse
René Diff, Diff
Bruger sans serif
René Diff, René Diff
Farvel nu!
Nå øv du

\endverse
\endsong



\beginsong{Våbenfysik}[sr={Melodi: Jutlandia}%37
,
by={FysikRevy 2003}
,
cr={}]\hypertarget{Våbenfysik}{}
\label{song37}

\beginverse
Det var i 1945 og nu
ville man ha’ fred,
men der var krig i Japan.
Der blev samlet uran nok og det
blev kastet ned,
for der var krig i Japan.
De fik at se hvad fysikken formår,
og sjove børn de næste mange år.

\endverse
\beginverse
\emph{[Omkvæd:]
Hey-Ho for våbenfysik!
Vi blæser på alle traktater.
Bomber her, bomber der, bomber
for fred!
Hvad skal man med diplomaaaaater?}

\endverse
\beginverse
Vi flyver gennem natten og gi’r
din fjende smæk.
Han ser os ikke komme.
Vi stiller ingen spørgsmål, og
så snart vi har din check,
så er krigen omme.
Hvis du gi’r, fyrer vi den sgu af.
Det er det, de unge vil ha’.

\endverse
\beginverse
\emph{[Omkvæd]}

\endverse
\beginverse
Vores salgskontor har åbent hele
døgnet – bare ring;
du skal ikke tøve.
Vacuum, brint og EMP og andre
sjove ting,
dem vil vi gerne prøve.
Her er ugens tilbudskatalog:
Start tre krige og betal for to!

\endverse
\beginverse
\emph{[Omkvæd]}

\endverse
\beginverse
\emph{[Omkvæd (ændret):]
Hey-Ho for våbenfysik!
Vi blæser på alle traktater.
Bomber her, bomber der, bomber
for fred!
Vi nakker de slyngelstaaaaater.}

\endverse
\endsong



\beginsong{Terminalsangen}[sr={Melodi: Shubidua - vuffelivov}
,
by={DIKUrevy 1978}
,
cr={}]\hypertarget{Terminalsangen}{}
\label{song38}

\beginverse
Jeg har en skærm med mange taster
En for hvert symbol
Og bagved sidder lysintensiteten
Den ledning har mange tråde
En til hver sin bit
og en ekstra en til pariteten

\endverse
\beginverse
Når man har venner og kærester, så er man normal
Men de ta'r tiden fra mig og min terminal

\endverse
\beginverse
Jeg er koblet via DIXI
Når DIXI ellers vil
Og der er plads på centrets multiplekser
Når jeg har lyst så kan jeg sidder
Og lege natten lang
Med RECKUs mange programmelkomplekser

\endverse
\beginverse
Når man har venner og kærester så er man normal
Men de ta'r tiden fra mig og min terminal

\endverse
\beginverse
Jeg kør' på en maskine
Der klarer tusind jobs
Selvom deta'r syv lange og syv bredde
CAU'en har den to af
Og det er vældig smart
En til hvis den anden sku' vær' nede

\endverse
\beginverse
Når man har venner og kærester så er man normal
Men de ta'r tiden fra mig og min terminal

\endverse
\beginverse
Jeg spiller skak og kryds og bolle
Den hele lange nat
Det er nu trist man ingen kender
For selvom den er dejlig
Så er den datamat
Nu kun et surrogat for menn'ske-venner

\endverse
\beginverse
Når man har venner og kærester så er man normal
Og har det bedre end mig med min terminal

\endverse
\endsong



\beginsong{Spas i Waders}[sr={Melodi:Hit'n'Hide - Space Invaders}
,
by={Biorevy 2012}
,
cr={}]\hypertarget{Spas i Waders}{}
\label{song39}

\beginverse
Spas i waders i en så
Fanger Myggelarver, Super biolog
Spas I waders I en så
Se på makrofytter, mål pH og lys

\endverse
\beginverse
Må afsted, kom nu med

\endverse
\beginverse
Du er feltbiolog med super seje waders på
Render og ta'r prøver dagen lang
Jeg er feltbiolog med super seje waders på
Renser danske søer for plankton

\endverse
\beginverse
Spas i waders i en sø ...

\endverse
\beginverse
Spas i waders i en sø...

\endverse
\beginverse
Må afsted, kom nu med

\endverse
\beginverse
Du er feltbiolog udstyret med planktonnet
Tæller og nøgler fisk og zooplankton
Ude i søerne, der har jeg net med småfisk i
At nøgle og måle, det kan jeg li'

\endverse
\beginverse
Spas i waders i en sø ...

\endverse
\beginverse
Spas i waders i en sø ...

\endverse
\beginverse
Du ta'r afsted op til Hillerød
Jeg vil med, se på livet i søen
 -- Vi finder en ny plankton art
Du og jeg tilsammen kan vi alt
 -- Vi kan alt

\endverse
\beginverse
Spas i waders i en sø ...

\endverse
\beginverse
Spas i waders i en sø ...

\endverse
\beginverse
Kom nu med!

\endverse
\endsong



\beginsong{Finite Simple Group (of order two)}[sr={}
,
by={The Klein Four Group}
,
cr={}]\hypertarget{Finite Simple Group (of order two)}{}
\label{song40}

\beginverse
The path of love is never smooth
But mine's continuous for you
You're the upper bound in the chains of my heart
You're my Axiom of Choice, you know it's true

\endverse
\beginverse
But lately our relation's not so well-defined
And I just can't function without you
I'll prove my proposition and I'm sure you'll find
We're a finite simple group of order two

\endverse
\beginverse
I'm losing my identity
I'm getting tensor every day
And without loss of generality
I will assume that you feel the same way

\endverse
\beginverse
Since every time I see you, you just quotient out
The faithful image that I map into
But when we're one-to-one you'll see what I'm about
'Cause we're a finite simple group of order two

\endverse
\beginverse
Our equivalence was stable,
A principal love bundle sitting deep inside
But then you drove a wedge between our two-forms
Now everything is so complexified

\endverse
\beginverse
When we first met, we simply connected
My heart was open but too dense
Our system was already directed
To have a finite limit, in some sense

\endverse
\beginverse
I'm living in the kernel of a rank-one map
From my domain, its image looks so blue,
'Cause all I see are zeroes, it's a cruel trap
But we're a finite simple group of order two

\endverse
\beginverse
I'm not the smoothest operator in my class,
But we're a mirror pair, me and you,
So let's apply forgetful functors to the past
And be a finite simple group,
be a finite simple group,
Let's be a finite simple group of order two (Why not three?)

\endverse
\beginverse
I've proved my proposition now, as you can see,
So let's both be associative and free
And by corollary, this shows you and I to be
Purely inseparable. Q.E.D.

\endverse
\endsong



\beginsong{Livstræet}[sr={}
,
by={Erik Lindebjerg og Hans Holm}
,
cr={}]\hypertarget{Livstræet}{}
\label{song41}

\beginverse
Der er så meget, der kan trykke,
gøre dagen trist og grå.
Se de folk, der uden lykke
bare går og går i stå.
\emph{Lad dem lege i livstræets krone,
lad dem føle, at livet er stort,
lad dem skue de blå horisonter
og himmelhvælvingens port.}

\endverse
\beginverse
Der er så meget uden varme,
uden ånd og uden liv.
Folk bli`r fattige og arme,
tænker kun på tidsfordriv.
\emph{Lad dem lege\ldots}

\endverse
\beginverse
Der er så mange uden venner,
uden kærlighed og kys.
Folk, som ingen andre kender
hvor skal de mon finde lys?
\emph{Lad dem lege\ldots}

\endverse
\beginverse
Der er så hård en kamp om magten,
alle kæmper for sig selv.
Er der ingen, der har sagt dem,
at de slår sig selv ihjel?
\emph{Lad dem lege\ldots}

\endverse
\beginverse
Der er så mange, der er slaver,
lænket fast til job og tid.
Hvorfor bruges vore gaver
uden tanke, uden vid?
\emph{Lad dem lege\ldots}

\endverse
\endsong



\beginsong{I Morgen er Verden Vor}[by={},
sr={Melodi:Tomorrow belongs to me - Cabaret}
,
cr={}]\hypertarget{I Morgen er Verden Vor}{}
\label{song42}

\beginverse
Se solen, der skinner
på kalv og på kid.
Se parken, der dufter af vår.
Nu sammen vi hilser
den nye tid.
I morgen er verden vor!
\endverse
\beginverse
Professoren giver
af alt hvad han ved.
Studerende sandheden får.
Og æren den venter
på os et sted.
I morgen er verden vor!
\endverse
\beginverse
Der tastes og kigges
på skærmene - se!
Det spirer, hvor Haarder han sår.
Men snart hvisker alle
BRUG  E D B
I MORGEN ER VERDEN VOR!
\endverse
\beginverse
ÅH EDB, EDB,
DU ER DET ORD,
DER FØRER OS FREM MOD DE ÅR,
HVOR VERDEN BLI'R STYRET AF DATALOG'R
I MORGEN ER VOR
I MORGEN ER VOR
I MORGEN ER VERDEN VOR!!
\endverse
\beginverse
\tiny{Gå direkte til 43!}
\endverse
\endsong


\beginsong{DAT62(1/2)80 Slagsang}[by={},
sr={Melodi:Fy Fy Skamme Skamme - McEwan og Ingerslev}
,
cr={}]\hypertarget{DAT62(1/2)80 Slagsang}{}
\label{song43}

\beginverse

\endverse
\beginverse
Må vi spy i datamaten,
må vi tæve den med køller
må vi smadre automaten
for at få en masse øller
må vi gå til forelæsning
i DIKUs høje sale
må vi råbe øv og foo og bar
og håndkøre som gale

\endverse
\beginverse
Ja ja ja ja ja - det må vi gerne

\endverse
\beginverse
Ørhn Ørhn hent hent skål skål spy spy
toogtredshalvfirs hent skål skål spy spy
- det var flot!

\endverse
\beginverse
Ja det må vi godt!

\endverse
\endsong


\beginsong{Bobbelsortering}[sr={Melodi:Liva Weel - Lad det boble}
,
by={DIKUrevy 2013}
,
cr={}]\hypertarget{Bobbelsortering}{}
\label{song44}

\beginverse
Ret din peger mod dit datasæt
følg nu med, lad vær' at blive træt.
Algoritmen listen traverserer
data vi sorterer ret.

\endverse
\beginverse
Tag to tal og se på deres værdi
er det mindre, bytter vi fordi
facit søges, korrekthed øges.
Så'n er datalogi.

\endverse
\beginverse
Lad dem boble, boble, boble,
bytte om på de her to
blot programmet terminerer,
kan jeg sagtens her bestå.

\endverse
\beginverse
Når man først sit emne trækker
og hukommelsen den strejker
håber gavmildheden rækker
bare manglerne er få.

\endverse
\beginverse
Jeg behøver aldrig lære fler'
hvis jeg kan bevise denne her
andre algoritmer de er svære
siger vores lærebog.

\endverse
\beginverse
Når engang jeg bobler helt på plads
er al tiden gået er mit sats.
Min metode er langsom kod'
Pawels spørgsmål gør nas.

\endverse
\beginverse
Lad dem boble, boble, boble,
bytte om på de her to
vi kan se det terminerer,
ergo pensum jeg forstår.

\endverse
\beginverse
Når I nu om lidt voterer
skal I husk' hvad jeg præsterer:
Algoritmen terminerer
så lad mig nu bestå!

\endverse
\endsong



\beginsong{Målteorisangen}[sr={Melodi: Vi lister os afsted på tå}
,
by={Matematikrevyen 2006}
,
cr={}]\hypertarget{Målteorisangen}{}
\label{song45}

\beginverse
Vi lister os af sted på tå
når vi skal integrere.
Vi måler kun på det, vi må,
alt andet la’r vi være.
Jeg måler på en gammel spand.
Jeg måler på din tissemand.
Og ellers måler vi det, vi nu kan,
både Tage og Flemming og Christian!

\endverse
\beginverse
Vi starter med et Radon-mål
og fjerner dig fra støtten,
og pluds’lig har du mål på nul
hvor før du havde sytten.
Du tro’de, at du var no’ed stort,
nu er du bar’ en lille lort,
men det må vi gerne fordi vi har spurgt
både Flemming og Tage og Christian!

\endverse
\beginverse
Og Kurzweil-Henstock integralet
det er helt umuligt!
Halvtreds procent det rene pral,
halvtreds procent ubru’ligt!
Og Riemann han blev pluds’lig væk,
måske var det en fejl i \TeX,
men så kan vi bruge den kære Lebesgue,
både Flemming og Tage og Christian!

\endverse
\beginverse
[Siges:] \emph{Arjjjj! Du må ikke måle på min tissemand!}
\emph{Bare rolig, jeg bruger tællemålet.}
Men hov, det kan jo ikke gå!
Den her kan du ikke måle på!
For udvalgsaksiomet det stoler vi på,
både Tage og Flemming og Christian!

\endverse
\endsong



\beginsong{Internetdiskurs}[sr={Melodi: 4 Non Blonde - What's Up}
,
by={DIKUrevy 2014}
,
cr={}]\hypertarget{Internetdiskurs}{}
\label{song46}

\beginverse
Jeg har altid syn's at det var let
(At) finde spas på internet
og mail
og og-så på YouTube

\endverse
\beginverse
Der sker det tit   at jeg gør mig klar
til at skrive en   farlig kommentar
af had
Når an-dre taler pænt

\endverse
\beginverse
  Jeg kan at du har en veltrænet kat
Træner du den dog mon med tunsalat?
jeg gør, men nu er den, så tyk den kan rulle

\endverse
\beginverse
Nu går jeg frivilligt ud af mit gode skind
For nu
skal de sættes på plads
"Hør lige her!"

\endverse
\beginverse
"GAYAYAY"
"RON PAUL 2012!"

\endverse
\beginverse
"GAYAYAYAYAY"
"IK(KE) RACIST, MEN..-"

\endverse
\beginverse
Mmmmh..
Mmmmh, mmmhmmmhm(SVIN!)mmhmmm
Mmmhmmmhmmmhmmm

\endverse
\beginverse
Det er hå-årdt
Åh min Gud, det er hå-årdt
For folk de tager fejl
På internettet!!

Alle de folk burd' bannes  fra   internet
Men de plejer nu mest at banne mig
Åh nej
  Hvor skal jeg, hvor kan jeg nu skrive

\endverse
\beginverse
Jeg tror jeg prø-ver med /b/ inde på fire-chan
Men jeg bli'r skræmt af en blo-dig tissemand
Den er
 Så klam at jeg bli'r bang'
"Lad vær' med det!"

\endverse
\beginverse
"GAYAYAY"
"FUCK YOU AND DIE!"

\endverse
\beginverse
"GAYAYAY"

\endverse
\beginverse
"GAYAYAY!"

\endverse
\beginverse
giver til sidst op i afmagt

\endverse
\beginverse
Nu er der kun én side som er rar
og man får karma for at vær' en nar
åh ja...
Hurra for Reddit

\endverse
\endsong



\beginsong{Biologiske Våben}[by={MBKrevy 2012}
,
sr={Melodi: Nordahl Grieg - Kringsav av fiender}
,
cr={}]\hypertarget{Biologiske Våben}{}
\label{song47}

\beginverse
Genmodifikation
Gør mig til gud!
Skaber en pathogen
Sender den ud!
Hvad vil du gøre nu?
Din krop er åben:
Sygdomment venter dig,
Den er mit våben

\endverse
\beginverse
Giftstoffer spreder sig,
Gennem din krop:
Skabt til at ramme alt,
Gi'r hjertestop.
Lammer først nerverne,
Kvæler din råben,
Toksiner skader dig
De er mit våben!

\endverse
\beginverse
Mød mutationerne
Du står for skud,
Tumorne vokser frem,
Under din hud!
Uhæmmet cellevækst,
Skrigen og råben
Kræften den spreder sig
Den er mit våben!

\endverse
\beginverse
Din krop, min legeplads,
Vi gi'r dig HIV!
Dør dit immunforsvar,
Ta'r vi dit liv
Du får en infektion
Drop blot din håben
For retrovirussen
Den er mit våben.

\endverse
\beginverse
Løftet det lyder nu
Forsker til folk
Vi vil jer skade med
Forskningens dolk
Dit barn vil fødes dødt
Hør mødres råben
Angsten vil sprede sig
Mærk vores våben!

\endverse
\endsong



\beginsong{Kvanter i måneskin}[sr={Melodi: Danse i måneskin}
,
by={FysikRevy 2004}
,
cr={}]\hypertarget{Kvanter i måneskin}{}
\label{song48}

\beginverse
Ved tavlen er jeg en charmør,
og pigerne bli’r bløde som smør.
Når jeg pertuberer, mere,
så får jeg topkarakterer.
I dag er det fredag, jeg går
ned på Cafeen? og får
en Guld Tuborg fra hanen, til ganen.
Og så går jeg hjem til mig selv.

\endverse
\beginverse
%[Omkvæd:]
Og kvanter i måneskin.
Jeg regner ud og ind
på energi og spin.
Kvanter i måneskin.

\endverse
\beginverse
Når piger de synes, jeg er sær,
så si’r jeg: ”Det’ sådan jeg er.”
Jeg regner på atomer, tro mig,
der er mindst ti millioner.
Jeg sidder og læser Euklid,
Hun siger, jeg spilder min tid.
Hun siger: ”Jeg vil ha’ dig. Ta’ mig!”
Og så går jeg hjem til mig selv!

\endverse
\beginverse
%[Omkvæd]
\lrep Og kvanter i måneskin.
Jeg regner ud og ind
på energi og spin.
Kvanter i måneskin.\rrep

\endverse
\endsong



\beginsong{Doom Doom}[sr={Melodi: Mabel - Boom Boom}
,
by={DIKUrevy 2006}
,
cr={}]\hypertarget{Doom Doom}{}
\label{song49}

\beginverse
Hvad er det der galt med mig? Jeg føler mig lidt sløj.
Jeg har fået min arm skudt af, og det er noget møg.
Der løber blod i øjet, mine tarme falder ud.
Det er vist ikke helt så godt hvis jeg ta'r flere skud.

\endverse
\beginverse
Nej, jeg gik op til lægen; og ved du hvad han sag'?
``Hvis du skal lege superhelt, så skal du ikke klag'.
Så tag dig her en tudekiks, tør øjnene og skær.
Jeg ved jo at din motorsav den står dig meget nær.''

\endverse
\beginverse
Det er helt naturligt, blodet sprøjter løs,
når en super shotgun er din ven.
(men) Har man ingen hitpoints,
så bli'r man let nervøs.
Så tag dig blot et medikit, og så' du frisk igen.

\endverse
\beginverse
Jeg spiller Doom, Doom for Doom er min trøst,
så længe hjertet banker her i mit bryst.
Jeg spiller Doom, Doom og spiller som besat.
Jeg lemlæster og myrder dag og nat.

\endverse
\beginverse
Er det den sidste bane? Jeg fatter ik' et quake.
Nu ser jeg rulleteksterne, monstrene er væk.
Fuck! Jeg har gennemført det. Hvad skal jeg ta' mig til?
Mit liv er fand'me kedeligt -- et pissedårligt spil.

\endverse
\beginverse
Det er helt naturligt, blodet sprøjter løs,
når en super shotgun er din ven.
(men) Har man ingen hitpoints,
så bli'r man let nervøs.
Så tag dig blot et medikit, og så' du frisk igen.

\endverse
\beginverse
Jeg spiller Doom, Doom og Doom Doom igen.
Min BFG har altid vær't min bedste ven.
Jeg spiller Doom, Doom og spiller som besat.
Jeg lemlæster og skyder dag og nat.

\endverse
\beginverse
Jeg spiller Doom, Doom og Doom Doom igen.
Min BFG har altid vær't min bedste ven.
Jeg spiller Doom, Doom og spiller som besat.
Jeg lemlæster og skyder dag og nat.

\endverse
\beginverse
Jeg spiller Doom, Doom og spiller som besat.
Jeg lemlæster og skyder dag og nat.

\endverse
\endsong



\beginsong{Humanist}[sr={Melodi:Dario von Slutty - Onani}
,
by={DIKUrevy 2013}
,
cr={}]\hypertarget{Humanist}{}
\label{song50}

\beginverse
Humanist - en KUAist
Humanist - en KUAist
Humanist - det' godt nok trist:
Det ta'r dem 10 semestre
og kan ikke brug's til sidst.

\endverse
\beginverse
Nu skal I hør':
Jeg er ham der Klaes, jeg har ECTS-behov,
men kurserne jeg tager, de gi'r mig aldrig lov.
De taler mig i søvne, her til forelæsningen
Og når den så er ovre, har jeg det hele glemt igen.

\endverse
\beginverse
Hvorfor tage kurser, der altid dumper mig?
Det er derfor, at jeg siger: "Der er en nem're vej!"

\endverse
\beginverse
Humanist - en KUAist ...

\endverse
\beginverse
Humanist - en KUAist ...

\endverse
\beginverse
Se, de fleste går og drømmer om at arbejde med kode,
men jeg vil bar' være færdig - arbejdsløshed er på mode!
Jeg er ligeglad med penge, jeg har børneopsparing
Og når den så løber tør, kan jeg pizza udbring'.
Hvorfor skulle arbejd', og ikke ha' tid til leg?
Det er derfor, at jeg siger: "Der er en nem're vej"

\endverse
\beginverse
Humanist - en KUAist ...

\endverse
\beginverse
Humanist - en KUAist ...

\endverse
\beginverse
Nu er det ovre, jeg er endelig blevet fri
Er blevet kandidat i østrigsk eskimologi
Der findes intet job hvori denne grad kan bruges
Så alle jobsamtaler sluttes af med at der buh'es.

\endverse
\beginverse
Her er det så, jeg fortryder jeg droppede ud
Hvis bare jeg kunne kode, så fremtiden lysere ud

\endverse
\beginverse
Humanist - en KUAist ...

\endverse
\beginverse
Humanist - en KUAist ...

\endverse
\beginverse
Datalog - han er jo god
Datalog - han er jo god
Datalog - får job, værsgo'
Det ta'r dem 20 semestre
men i det mindste kan de kode.

\endverse
\beginverse
Datalog - han er jo god ...

\endverse
\beginverse
Det ta'r dem 20 semestre
men i det mindste kan de kode.

\endverse
\endsong



\beginsong{Doktrinen}[sr={Melodi:Otto Brandenburg - Alle sømænd er glade for piger}
,
by={DIKUrevy 2009}
,
cr={}]\hypertarget{Doktrinen}{}
\label{song51}

\beginverse
Alle russer skal lære doktrinen
man skal kode i Moscow ML
men de modsætter sig disciplinen
når de sidder bag skærmene selv

\endverse
\beginverse
Og koder i Java, C\# eller Ada
selv Fortran og Matlab, og HTML
Foruden lidt Delphi, og Python og Ruby
og Visual Basic med MySQL

\endverse
\beginverse
Som instruktor er ansvaret mit
russers uvaner skal rettes ind
de skal ik' tro at valget er frit
for de lærer jo ikke en pind

\endverse
\beginverse
Når de koder i Java, C\# eller Ada ...

\endverse
\beginverse
Efter tolv år er jeg kandidat
Kapitalen den kalder på mig
min doktrin havde nul resultat
for på arbejdet der sidder jeg

\endverse
\beginverse
Og koder i Java, C\# eller Ada ...

\endverse
\endsong



\beginsong{\TeX'e-sangen}[sr={Melodi: Sebastians Skatteøen - Ostesangen}
,
by={Matematikrevyen 2012}
,
cr={}]\hypertarget{TeX'e-sangen}{}
\label{song52}

\beginverse
Jeg skal rette opgaver fra RUSser.
Driver mig til vanvid,
intet står som jeg vil ha' det.
Alle mine russer går og sjusker,
og jeg bruger lang tid.
Maple bruger de til alt.
Bar' de ku brug':
\endverse

\beginverse
\texttt{fancyhdr}, \texttt{itemize},
\texttt{verbatim} og \texttt{normalsize},
hvis det skal vær' tjekket.
\texttt{tabular}, \texttt{$\backslash$begin$\{$align$\}$},
hundredvis af \texttt{\$}'tegn,
så skal det vær' \TeX'et.
\TeX, jeg vil ha' \TeX, jeg må ha' \TeX.
\TeX, jeg vil ha' \TeX, jeg må ha' \TeX.
\TeX, \TeX, \TeX, \TeX, \TeX, jeg må ha' T\!$_{\text{\normalsize E}}$-$_{\text{\normalsize EEEE}}$\!X.
\endverse

\beginverse
Mellemugen bringer os et kursus.
Fjerner al den vanvid;
jeg skal lære dem at \TeX'e.
Selvom det ik' er den rene luksus,
lønner det sig altid
når de taler \TeX'ens sprog.
Så kan de brug':
\endverse

\beginverse
\texttt{fancyhdr}, \texttt{itemize}, \ldots
\endverse

\beginverse
Nu bru'r de al tiden på at \TeX'e.
Driver dem til vanvid,
det får de snart mer' end nok af.
Nu er alle opgaverne lette
-- har så meget fritid,
får min nattesøvn
og drømmer nu så sødt om:
\endverse

\beginverse
\texttt{fancyhdr}, \texttt{itemize}, \ldots
\endverse

\beginverse
\texttt{fancyhdr}, \texttt{itemize}, \ldots
\endverse
\endsong



\beginsong{Hjemmehackeriet}%53
[sr={Melodi: Hjemmebrænderiet}
,
by={DIKUrevy 2001}
,
cr={}]\hypertarget{Hjemmehackeriet}{}
\label{song53}

\beginverse
Jeg bor her i Ishøj på syvende sal
i en lejlighed der stort set er normal.
En stue, et køkken, et bad med WC
og et kammer hvor jeg har min hjemme-PC.

\endverse
\beginverse
\emph{Jeg hacker, jeg cracker, jeg downloader spil,
og jeg logger ind lig' præcis hvor jeg vil.
Jeg kender dit password, jeg læser din post;
for en hacker som mig er den slags hverdagskost.}

\endverse
\beginverse
Min fætter har hacket i Pentagons net.
De tro'ed det var svært, men han syn's det var let.
De fandt ham dog efter en længere jagt,
så nu er han ansat som sikkerhedsvagt.

\endverse
\beginverse
\emph{Jeg hacker, jeg cracker\ldots}

\endverse
\beginverse
Jeg laved' en virus som hed ``I Love You''.
Jeg indrømmer dog, jeg fortryder det nu.
Da jeg gik i banken, min løn for at få,
havde virusen sat der's computer i stå.

\endverse
\beginverse
\emph{Jeg hacker, jeg cracker\ldots}

\endverse
\beginverse
Hvis du sku' få lyst til at hacke lidt selv,
jeg ønsker dig al mulig lykke og held.
Det giver dig magt som om du var en gud,
og du kan endda få din pizza bragt ud.

\endverse
\beginverse
\emph{Jeg hacker, jeg cracker\ldots}

\endverse
\endsong



\beginsong{Studiedag}%54
[sr={Melodi: Hjemmebrænderiet}
,
by={Matematikrevyen 2012}
,
cr={}]\hypertarget{Studiedag}{}
\label{song54}

\beginverse
Min time den starter ved kvart over ni.
Jeg stiller min cykel, når klokken slår ti.
Jeg skynder min ned for at lære lidt SaSt,
og jeg håber at lidt af det sætter sig fast.
\endverse

\beginverse
\emph{Jeg elsker at læse på matematik.
De andre er rare, og så bli'r jeg kvik.
De regler om spredning dem kender jeg vagt,
men jeg håber da stadig jeg bli'r kandidat.}
\endverse

\beginverse
En gang efter frokost der har vi LinAlg.
En matrix er vigtig, har Morten fortalt.
Jeg bli'r lidt forsinket, det bli'r jeg konstant,
og jeg fatter vist aldrig den determinant.
\endverse

\beginverse
\emph{Jeg elsker at læse på matematik.
De andre er rare, og så bli'r jeg kvik.
De regler om søjler dem kender jeg vagt,
men jeg håber da stadig jeg bli'r kandidat.}
\endverse

\beginverse
Jeg ta'r på Caféen? og køber et shot,
men dankort-maskinen den afviser blot.
Mit kort er vist spærret -- jeg fatter det ik',
for min SU bør gå ind per automatik.
\endverse

\beginverse
\emph{Jeg elsker at læse på matematik.
De andre er rare, og så bli'r jeg kvik.
De regler om SU dem kender jeg vagt,
men jeg håber da stadig jeg bli'r kandidat.}
\endverse

\beginverse
Jeg dumper LinAlg og sidenhen SaSt,
så dropper jeg studiet og kører helt fast.
Jeg skylder en hel masse penge til mor,
og det er et problem for hun ved hvor jeg bor.
\endverse

\beginverse
\emph{Jeg elsker at læse på matematik.
De andre er rare, og så bli'r jeg kvik.
Men regler i Netto dem kan jeg forstå,
for jeg ved båd' hvor ketchup og dåser skal stå.}
\endverse
\endsong



\beginsong{Hjemmebrænderiet}%55
[sr={}
,
by={De Gyldne Løver -- Variant: TÅGEKAMMERET}
,
cr={}]\hypertarget{Hjemmebrænderiet}{}
\label{song55}

\beginverse
Jeg bor her i Stenhøj i et skævt lille hus
alene rent bortset fra rotter og mus,
i gården et lokum og bagved et skur
og i det stod mit brændevinsapparatur.
Det syder af fusel og bobler af gær,
din lever ta'r skade og øjet får stær,
det smager som rævepis og terpentin,
men det er billigt og så bli'r man fuld som et svin.
\endverse

\beginverse
Min bedstemor brændte sin brændevin selv
det slog både hende og manden ihjel,
de drak aldrig selv, men for heden den dag
brænderiet røg i luften med et helvedes brag.
Von Vusel geschnurrt und von Gehre gestört,
dein Liebchen macht scade und augen geschlört.
Von Fuchspiss geschmächt, und von Scheisse geschoffen,
Es ist billig und du wirst so wunderlich soffen.
\endverse

\beginverse
Min fader han brændte i $67$ år,
men så blev han snuppet af tolden i går,
så moar har skaffet et nyt apparat
og sat produktionen i gang i en fart.
It’s seething like fussel and boiling like gear,
your liver takes shade, and the eye may have stair.
It’s tasting like foxpis and may be some ass,
but it’s cheaper, and you will be drunk with a crash.
\endverse

\beginverse
Det hænder jeg selv ta'r en ordentlig søp,
mens jeg lytter til dråbernes sagte dryp, dryp.
Det er som musik fra et fint instrument,
for jeg ved det er sprit på omkring $60\%$.
Kamsjatka fulaska boblinska gaeruska,
dju lev skinavara na oeleski hi pip.
Sja smakal urinos Breznjevski uhumsk,
Sie lahanski virundski paagulski pladask.
\endverse

\beginverse
Og hvis du har får lyst til at smage det selv,
så kom kun til mig helst når dagen går på hæld,
slå tre sla’v på ruden og stik mig en ti’r,
så får du en sjat af min livseleksir.
El bruso fusellos boblina gaerror,
levino ituo, e cicco staesor.
El taste pissare e Costa del Sol,
nul monetos, e du havor ondt i din knold.
\endverse

\beginverse
Mig bo ær i stenhøj i sø lille us.
Alene maj sæl og et sø lille mus.
Om ba ve dar ær 2iletter mæ brat.
Og dær maj har gæmt mine spritapparat.
Mig gæmmer det gåt, for policen ka li.
Og maj ta i nakken ov sæn te Tyrki.
Å det maj vær majet stor ked a fordi.
Måske maj i fængelse stor lang ti skab li.
\endverse
\endsong



\beginsong{Sensommervisen}[sr={}
,
by={Kirsten og Finn Jørgensen}
,
cr={}]\hypertarget{Sensommervisen}{}
\label{song56}

\beginverse
Æbler lyser rødt på træernes grene,
høsten går ind.
Går igennem skoven ganske alene,
stille i sind.
\emph{Gyldne farver og sensommerbrise
fylder hjertet med vemodig musik
Går og nynner en sensommervise
fjernt fra byens larmende trafik.}
\endverse

\beginverse
Sommerbrisen danner krusning på søen,
mystisk og sort.
Stæreflokke svæver højt over øen
snart ta'r de bort.
\emph{Gyldne farver\ldots}
\endverse

\beginverse
Duft af brænderøg blandt brunlige bregner,
blåsorte bær.
Stille summen mellem blade, som blegner
aft'nen er nær.
\emph{Gyldne farver\ldots}
\endverse

\beginverse
Modne rønnebær bag dybgrønne grene
rødt titter frem.
Går igennem skoven ganske alene
nu må jeg hjem!
\emph{Gyldne farver\ldots}
\endverse
\endsong



\beginsong{Sortsnak}[sr={Melodi: Cæsar -- Storkespringvandet}
,
by={}
,
cr={}]\hypertarget{Sortsnak}{}
\label{song57}

\beginverse
\large
$\aleph\ \aleph\ \aleph\ \aleph\ \chi$
$\aleph\ \aleph\ \aleph\ \pi$
$\aleph\ \pi\ \chi\ \aleph\ \pi\ \chi\ \pi$
$\chi\ \chi\ \pi\ \aleph\ \lambda\ \aleph$

\endverse
\beginverse
[For de øvede:]
\large
$\beta\ \beta\ \beta\ \beta\ \xi$
$\beta\ \beta\ \beta\ \psi$
$\beta\ \psi\ \xi\ \beta\ \psi\ \xi\ \psi$
$\xi\ \xi\ \psi\ \beta\ \gamma\ \beta$

\endverse
\endsong


\beginsong{Så spiller vi spil}[sr={Melodi: Höhner - Viva Colonia}
,
by={DIKUrevy 2006}
,
cr={}]\hypertarget{Så spiller vi spil}{}
\label{song58}

\beginverse
For en fattig studer'ne
er dagligdagen svær.
Læser mange tykke bøg'r,
laver meget her.
Alle vore timer går med hjernegymnastik;
hvordan vi kan klare det, forstår man ik'.

\endverse
\beginverse
Vi har så travlt til dagligt,
der kræves mange svar.
Når opgaven er færdig,
så er den næste klar.
Vi får skrevet hundred' sider
i gennemsnit hver dag,
men når en weekend nærmer sig, ja så
så slapper vi af.

\endverse
\beginverse
\emph{[Omkvæd] Så spiller vi spil, hele dagen
ved datamaterne.
Vi lever og ånder
for Doom og Quake med mer',
med World of Warcrack morer vi os
her foran vor's PCer.}

\endverse
\beginverse
Dengang i gamle dage
da gjord' rapporter ondt.
Så fik vi blokstrukturen,
men det var ikke sundt.
Vi har stress og dårlig' nerver.
Vi sveder angstens sved.
Og beder tit en stille bøn til at
at DIKU går ned.

\endverse
\beginverse
\emph{[Omkvæd] Så spiller vi spil\ldots}
\emph{[Omkvæd] Så spiller vi spil\ldots}

\endverse
\beginverse
Og pluds'lig er det morgen,
så hurtigt kan det gå.
Man burde måske læse,
hvis nu vi skal beståååååååå....
{}[Siges] men det bliver ikke... nu.

\endverse
\beginverse
\emph{[Omkvæd, stille og langsomt stigende] Nu spiller vi spil\ldots}
\emph{[Omkvæd] Så spiller vi spil\ldots}

\endverse
\beginverse
Så spiller vi spil!

\endverse
\endsong



\beginsong{Der står en printer}
[sr={Melodi: Der Bor en Bager}
,
by={Matematikrevyen 2015}
,
cr={}]\hypertarget{Der står en printer}{}
\label{song59}

\beginverse
Der står en printer
Helt op' på fjerde
Den printer gratis
Hvis man er lærer
Den printer store
Den printer små
Den printer nogle
Med farver på
\endverse

\beginverse
For alle andre
Nu håbet splintrer
Vi havde 1000
Nu har vi mindre
Og har du penge
Så kan du få
Men har du ingen
Så må du gå
\endverse


\endsong



\beginsong{Jeg' en nørd}[sr={Melodi: tema fra Pokémon}
,
by={FysikRevy 2010}
,
cr={}]\hypertarget{Jeg' en nørd}{}
\label{song60}

\beginverse
AD\&D det er min leg,
og go det er min sport.
Nyt tøj, det bytter jeg
til sjældne magickort.

\endverse
\beginverse
World of Warcraft er så fedt.
Jeg spiller natten lang.
Back to the Future har jeg set,
for 85. gang!

\endverse
\beginverse
JEG’ EN NØRD!
Jeg’ stolt af det.
Jeg læser XKCD.

\endverse
\beginverse
JEG’ EN NØRD!
Skriiiiiiiver i \LaTeX
og jeg tænder på $dx$.

\endverse
\beginverse
JEG’ EN NØRD!
Nu er jeg ved
at finde en at nørde med,
som ka’ ta’ min uskyldighed.

\endverse
\beginverse
JEG’ EN NØRD!
Og jeg er stolt af det!
Yeaaahh…

\endverse
\beginverse
Jeg er level 10 i sværd,
så spil nu ikke smart.
I min bæltetaske er
der terninger parat.

\endverse
\beginverse
Sig at det skal være os
og mærk min energi.
Når jeg læser Hitchhikers’
Guide to the Galaxy.

\endverse
\beginverse
JEG’ EN NØRD!
Jeg’ stolt af det.
Jeg læser XKCD.

\endverse
\beginverse
JEG’ EN NØRD!
Skriiiiiver i \LaTeX
og jeg tænder på $dx$.

\endverse
\beginverse
JEG’ EN NØRD!
Nu er jeg ved
at finde en at nørde med,
som ka’ ta’ min uskyldighed.

\endverse
\beginverse
JEG’ EN NØRD!
Jeg er stolt af det.

\endverse
\beginverse
Jeg er stolt af det.
JEG’ EN NØRD!

\endverse
\endsong



\beginsong{Koder Hele Natten}[sr={Melodi:Rage Against The Machine - Killing In The Name Of}
,
by={DIKUrevy 2011}
,
cr={}]\hypertarget{Koder Hele Natten}{}
\label{song61}

\beginverse
De bruger magt, til at rykke os til HCØ.
De bruger magt, til at fjerne vores kantine.
De bruger magt, til at fjerne vores rusture.
Nu er det nok, vi sætter hårdt mod hårdt
Vi er anonyme, vi er DIKU!

\endverse
\beginverse
Koder hele natten!

\endverse
\beginverse
Der er dem der koder VB, det er dem uden et CV
Der er dem der koder Brainfuck, de kan ikke få det hårdt nok
Der er dem der koder Java, det' det en'ste de kan lave.
Der er dem der koder ML, det er os der gør det hel' selv!
Kod!

\endverse
\beginverse
Koder hele natten!
Koder hele natten!

\endverse
\beginverse
I skærmens lys vi skal kode
Vi skal besejre hvor klode
Ved DIKU danner vi pagten
Med cyberkrig får vi magten
I skærmens lys vi skal kode
Vi skal besejre hvor klode
Ved DIKU danner vi pagten
Med cyberkrig får vi magten
I skærmens lys vi skal kode
Vi skal besejre hvor klode
Ved DIKU danner vi pagten
Med cyberkrig får vi magten!

\endverse
\beginverse
Rejs jer op, tøm jeres kop, i fællesskab får vi DIKU op!
For all' skal vid', med dataslid, vi genvinder nu DIKUs storhedstid!
Rejs jer op, tøm jeres kop, i fællesskab får vi DIKU op!
For all' skal vid', med dataslid, vi genvinder nu DIKUs storhedstid!

\endverse
\beginverse
Der er dem der koder VB, det er dem uden et CV ...

\endverse
\beginverse
Koder hele natten!
Koder hele natten!

\endverse
\beginverse
I skærmens lys vi skal kode
Vi skal besejre hvor klode
Ved DIKU danner vi pagten
Med cyberkrig får vi magten
I skærmens lys vi skal kode (hver en bit vi vil se)
Vi skal besejre hvor klode (bestem' alt der skal ske)
Ved DIKU danner vi pagten (datalog'r når' de størst)
Med cyberkrig får vi magten (nettet slukker hvor tørst)
I skærmens lys vi skal kode (hver en bit vi vil se)
Vi skal besejre hvor klode (bestem' alt der skal ske)
Ved DIKU danner vi pagten (datalog'r når' de størst)
Med cyberkrig får vi magten

\endverse
\beginverse
Rejs jer op, tøm jeres kop, i fællesskab får vi DIKU op!
For all' skal vid', med dataslid, vi genvinder nu DIKUs storhedstid!
Rejs jer op, tøm jeres kop, i fællesskab får vi DIKU op!
For all' skal vid', med dataslid, vi genvinder nu DIKUs storhedstid!
Kod nu!

\endverse
\beginverse
Kod! Kod nu!

\endverse
\beginverse
DIKU, det er nu de skal lære
DIKU, vi vil kæmp' for din ære!
DIKU, det er nu de skal lære
DIKU, vi vil kæmp' for din ære!
Hack dem, det er nu i skal kode
Hack dem, vi besejrer vor klode
Hack dem, det er nu i skal kode
Hack dem, vi besejrer vor klode
Moderkoder!
Kod!

\endverse
\endsong



\beginsong{Q.E.D.}[sr={Melodi:Within Temptation - Our Solemn Hour}
,
by={FysikRevy  2010}
,
cr={}]\hypertarget{Q.E.D.}{}
\label{song62}

\beginverse
Quod erat demonstrandum!
Med disse ord skal sandhed skrives
Quod erat demonstrandum!
Uvidenheden må fordrives

\endverse
\beginverse
Quod erat demonstrandum!
Quod erat demonstrandum!
Quod erat demonstrandum!

\endverse
\beginverse
Mørket samler sig
selv i den mindste krog
er logik og logos på tilbagetog
Intelligent design
og anden overtro
videnskaben taber til en astrolog

\endverse
\beginverse
Så vi må handle snart før det er forbi
Slut med tavshed nu det tid at si'e

\endverse
\beginverse
Quod erat demonstrandum!
Med disse ord skal sandhed skrives
Quod erat demonstrandum!
Uvidenheden må fordrives
Quod erat demonstrandum!
Skal verden vindes må vi st den bi
Naturvidenskabs logik vil gør' dig fri!

\endverse
\beginverse
En ungdomsgeneration
's fornuft sat på auktion
Vor fremtid smides bort
Vor viden kom til kort
Fra Unge Mødre og
til Paradise Hotel
vores hjerneceller keder sig ihjel

\endverse
\beginverse
Idiotiens lænker har holdt dig ned'
så kom ind i kampen og syng med!

\endverse
\beginverse
Quod erat demonstrandum! Med disse ord ...

\endverse
\beginverse
Quod erat demonstrandum! Med disse ord ...

\endverse
\beginverse
Quod erat demonstrandum! Med disse ord ...

\endverse
\beginverse
Quod erat demonstrandum!
Med disse ord skal sandhed skrives (Sandhed skrives)
Quod erat demonstrandum!
Uvidenheden må fordrives (må fordrives)
Quod erat demonstrandum!
Skal verden vindes må vi stå den bi
Naturvidenskabs logik vil går' dig fri!

\endverse
\endsong



\beginsong{Hvem Sidder Der Ved Skærmen}%63
[sr={Melodi: Jens Vejmand}
,
by={DIKUrevy 1973}
,
cr={}]\hypertarget{Hvem Sidder Der Ved Skærmen}{}
\label{song63}

\beginverse
Hvem sidder der ved skærmen
med strimler om sin hånd
med hulkort op til halsen
og om sin sko et bånd?
Det er såmænd en dat'er
som sidder dagen lang
ved skærmen og forvandler
sit Algol 5 til Slang.
\endverse

\beginverse
Og vågner du en morgen
i allerførste gry
og ser at skærmen skriver
på ny, på ny, på ny --
det er såmænd vor dat'er
som på dig sindssygt glor
og taster vildt og blodigt
ved terminalens bord.
\endverse

\beginverse
Og vandrer du til stuen
hvor alt er fint i stand,
og møder du en stakkel
hvis øjne står i vand --
det er såmænd vor Dat'er
som kommer dig på tværs
og ikke mer' kan finde
en fejl i line halvfjerds.
\endverse

\beginverse
Og vender du tilbage
i byger og i blæst,
mens aftenstjernen skælver
af kulde fra nordvest,
og klinger terminalen
bag ryggen ganske nær --
det er såmænd vor dat'er
som endnu sidder dér.
\endverse

\beginverse
Dog dataskærmen jævned
for ham den slemme vej,
men da det led mod julen,
da sagde skærmen nej;
det var såmænd vor dat'er
hans fingre slap den brat,
de bar ham væk fra stuen
en kold december-nat.
\endverse

\beginverse
Der er på opslagstavlen
et gammelt grønnet kort
hvor hulningen er ussel
og skriften ganske bort.
Det tilhørte vor dat'er
som aldrig gav et gny,
men på hans plads ved skærmen
der sidder nu en ny.
\endverse
\endsong



\beginsong{Hvem sidder foran skærmen}%64
[sr={Melodi: Jens Vejmand}
,
by={FysikRevy 2001}
,
cr={}]\hypertarget{Hvem sidder foran skærmen}{}
\label{song64}

\beginverse
Hvem sidder foran skærmen
med Cola i sin hånd,
med blege svage arme
og kroppen på et ton?
\emph{Det er jo datalogen,
en nørd, det er jo klart,
med koffein i blodet
og cybersex parat.}

\endverse
\beginverse
Og da den næste morgen
en kvinde træder ind
og vækker datalogen,
han fatter ik' en pind.
\emph{Det er hans lillesøster,
hun mangler sin BH;
og datalogen ryster,
det' ham der har den på.}

\endverse
\endsong



\beginsong{Jeg er fysiker}[sr={Melodi: Monty Python - The Lumberjack Song}
,
by={Matematikrevyen 2007}
,
cr={}]\hypertarget{Jeg er fysiker}{}
\label{song65}

\beginverse
Jeg er fysiker og jeg er glad.
Jeg kan min Schaums bog uden ad!
\endverse

\beginverse
Han er fysiker og han er glad.
Han kan sin Schaums bog uden ad!
\endverse

\beginverse
Ampere her, promille der,
jeg gør et Maple-plot.
Jeg plotter i den farve,
min kone si’r er flot!
\endverse

\beginverse
Ampere her, \ldots
\endverse

\beginverse
Han er fysiker \ldots
\endverse

\beginverse
Eksperimentér, analyser,
og spis en masse slik!
Jeg læser en artikel,
og leger med min rubikskube\ldots
\endverse

\beginverse
Eksperimentér, \ldots
\endverse

\beginverse
Jeg er fysiker, og jeg er fin.
Jeg kan formlen for min urin!
\endverse

\beginverse
Jeg rækkeudvikler på alt
fordi det har jeg lært.
Hvis blot jeg læste mat'matik
som ham der Erik Kjær!
\endverse

\beginverse
Han rækkeudvikler \ldots
\endverse
\endsong



\beginsong{$\alpha\beta$-sangen}[sr={Melodi: I en kælder sort som kul}
,
by={Tenna Schaldemose og Tine Nyeng, Majrevy 2004}
,
cr={}]\hypertarget{$alphabeta$-sangen}{}
\label{song66}

\beginverse
$\alpha\ \beta\ \gamma\ \delta$\hspace{-2.5pt}\protect\colorbox{white}{\phantom{$\delta$}}
\hspace{-2.5pt}$\delta$\hspace{-1.25em}\colorbox{white}{\phantom{$\delta$}}$\ \varepsilon\ \zeta$
$\eta\ \theta\ \iota\ \kappa$\hspace{-2.5pt}\protect\colorbox{white}{\phantom{$\kappa$}}
\hspace{-3pt}$\kappa$\hspace{-1.3em}\colorbox{white}{\phantom{$\kappa$}}$\ \lambda\ \mu\ \nu\ \xi$
$o\ \pi\ \rho\ \sigma$
$\tau\ \upsilon\ \varphi\ \chi\ \psi$
$\omega$ er sidste
$24$ på liste!
\endverse

\beginverse
Alfa beta gamma del-
ta epsilon zeta
eta theta iota kap-
pa lambda my ny ksi
omikron pi rho sigma
tau ypsilon phi chi psi
omega er sidste
fireogtyve på liste!
\endverse
\endsong



\beginsong{Åh, Maple}[sr={Melodi: Westlife - Mandy}
,
by={FysikRevy 2015}
,
cr={}]\hypertarget{Åh, Maple}{}
\label{song67}

\beginverse
Husker dengang jeg var rus
MatIntro, tog det lidt på sjus'

\endverse
\beginverse
En skygge af et blad
Et håb om at ku løse
Funktionen foran mig
Var det muligt

\endverse
\beginverse
Onsdag formiddag og jeg
Mærker "deadline" nærme sig
Skriver den funktion

\endverse
\beginverse
Og beder dig plot' den
Og alt du giver mig
Er blot en fejlmelding åh Maple

\endverse
\beginverse
For du gør mig så skingrende sindssyg
Når jeg plotter funktioner, ååh Maple
Selvom jeg prøver at være kærlig
Ignorerer du mig, åååh Maple

\endverse
\beginverse
Jeg sidder til eksamen nu
Er maple noget jeg ka' brug?
Men RootOfCSC?
Hvad fanden betyder det?
Har du andre svar
End errorfunktioner, Åh Maple

\endverse
\beginverse
For du gør mig så skingrende sindssyg
Når jeg plotter funktioner, ååh Maple
Når jeg bruger græske symboler
Ja, så fatter du hat! Åh Maple

\endverse
\beginverse
Du siger "export to Latex"
Hvad fuck er det, det virker ik'
Men jeg klare'd den eksam'
Jeg fik det hele gjort
Og hjælpen du gav mig
Var fem liniers rødt lort, åh Maple

\endverse
\beginverse
Bare tænk, hvis jeg var til eksamen
Uden dig på min PC, åh maple
Selvom jeg prøver at være kærlig,
Ignorerer du mig, åh Maple

\endverse
\beginverse
TI-89
Og Matlab og Mathcat
Er ubrugelig'
Men intet er værre end Maple

\endverse
\beginverse
For du gør mig så skingrende sindssyg
Når jeg plotter funktioner, ååh Maple
Selvom jeg prøver at være kærlig
Ignorerer du mig, åååh Maple

\endverse
\beginverse
Bare tænk, hvis jeg var til eksamen
Uden dig på min PC, åh maple
Selvom jeg prøver at være kærlig,
Skider du i munden på mig, åh Maple

\endverse
\endsong



\beginsong{Sangen om Nick}[sr={Melodi: Quangs sang}
,
by={}
,
cr={}]\hypertarget{Sangen om Nick}{}
\label{song68}

\beginverse
Her er en historie om en fyr der hedder Nick.
Han er tyve, bor på Ama'r, og han er ikke særligt kvik,
for han læser ude på KUA, han er nemlig humanist.
Og der er så mange piger -- hvorfor er han så så trist?
\endverse

\beginverse
Og Nick han mener ikke at han er en doven rad,
for han bruger hele dagen på at lære sine kvad.
Og så går han til eksamen, men det går altid skidt.
Så Nick han ville ønske at han var ham Kasper Smith.
\endverse

\beginverse
\emph{Så hva' fanden får dig til at sige du ikke er tilfreds,
at mat1GA det stinker, og dine lektier gi'r dig stress?
Du har alt hvad du skal bruge, Søren Eilers han er sej.
Der er tusind folk på KUA der ville ønske de var dig!}
\endverse

\beginverse
Nick har ingen far mer' og hans mor er meget syg,
så Nick må tjene penge inde i den store by.
Ja, Nick han har et studiejob, man sku' tro at det var fake,
for han oversætter tekster for TV3 til Ricki Lake.
\endverse

\beginverse
\emph{Så hva' fanden får dig til at tro at jeg vil se dig sidde og glan'?
Du sku' hell're ta' at bevæge dig ud i den komplekse plan.
Du har vektorrum og Cauchy-Schwarz, Borel og Weierstra\ss.
Husk på at ude på KUA der har de virk'lig keinen spa\ss!}
\endverse

\beginverse
Og når det så er fredag, er der ingen fredagsbar.
For han skal spare sine kræfter for at blive klar
til at skrive stil om ferien, hvor der var trange kår
med sin kæreste Anders Gaarde, som er 21 år.
\endverse

\beginverse
\emph{Så hva' fanden får dig til at tro at jeg vil høre på dit brok:
at du ikk' ka' li' din logik og ikk' får regnestykker nok.
Du har Dilans og Caféen?, du har tid til hurlumhej.
Der er tusind folk på KUA der ville ønske de var dig!}
\endverse

\beginverse
En sesquilinearform den ved han intet om,
og det er det der gør hans tilvær'ls' ååhh noget så tom.
For Nick han ville så gerne have lært om noget stort,
men det eneste han hører om er humaniora-lort!
\endverse

\beginverse
\emph{Så hva' fanden får dig til at sige at du har et dårligt liv?
Er du sød at prøve at se det i et større perspektiv?
Tag og tænk på andre end dig selv, tag ved lære af en som Nick.
Husk på du har det godt, når bare du læser mat'matik!}
\endverse
\endsong



\beginsong{Hvad skal vi \TeX'e i nat?}[sr={Melodi: Hvor skal vi sove i nat?}
,
by={DIKUrevy 2007}
,
cr={}]\hypertarget{Hvad skal vi TeX'e i nat?}{}
\label{song69}

\beginverse
Vi fik rapporten, og tænkte straks det samme;
går den, så går den, og drak os bankelamme.
Vi' alt for kloge, vi' nogle seje gutter.
Selv i en tåge, ta'r den kun fem minutter.

\endverse
\beginverse
Så fandt vi siden, hvor alting blev beskrevet.
Vi mangled' viden, for at få alt det lavet.
Vi hader livet under rapportregimer.
Vi er fortvivlet, vi har kun fjorten timer.

\endverse
\beginverse
Hvad skal vi \TeX'e i nat?
Vi har jo intet lavet.
Hvad skal vi \TeX'e i nat?
Vi er slet ikke færdig'.
Vi er jo fuldstændig sat.

\endverse
\beginverse
Vi emmed' af arrogancen og mistede chancen
for at gøre det godt.

\endverse
\beginverse
Hvad skal vi \TeX'e i nat?
Vi ku' jo lave genbrug.
Hvad skal vi \TeX'e i nat?
Så printer jeg den ud nu.
Vi ku' jo drikke en sjat.

\endverse
\beginverse
Vi ta'r en overspringshandling, det gi'r en forvandling
så vi tror det er flot.

\endverse
\beginverse
Vi lærte lektien, vi læser alle sider.
Så nu er vi mænd der ikke går og lider.
Vi starter tidligt og knokler meget længe.
Vi får lidt slidgigt -- ser sjældent vore senge.

\endverse
\beginverse
Men hov hvad er det, er det det gale pensum?
Ja, så forstår jeg der var så tavst på forum.
Vi starter forfra, vi føler os som tåber,
og i kantinen er det kun os der råber.

\endverse
\beginverse
Hvad skal vi \TeX'e i nat?
Vi har jo intet lavet.
Hvad skal vi \TeX'e i nat?
Vi er slet ikke færdig'.
Vi er jo fuldstændig sat.

\endverse
\beginverse
Vi emmed' af arrogancen og mistede chancen
for at gøre det godt.

\endverse
\beginverse
Hvad skal vi \TeX'e i nat?
Vi ku' jo lave genbrug.
Hvad skal vi \TeX'e i nat?
Så printer jeg den ud nu.
Vi ku' jo drikke en sjat.

\endverse
\beginverse
Vi ta'r en overspringshandling, det gi'r en forvandling
så vi tror det er flot.

\endverse
\endsong



\beginsong{Lad det gro}%70
[sr={Melodi:Frozen - Let It Go}
,
by={DIKUrevy 2014}
,
cr={}]\hypertarget{Lad det gro}{}
\label{song70}

\beginverse
Kantinen er tom, der er ingen at se
Og jeg drømmer for mig selv
Mit ønske er helt umuligt, (men) jeg har håb a-lligevel

\endverse
\beginverse
Men det vil ik' gro

\endverse
\beginverse
Det er det alle mennesker vil sig'
Gi' mig et skæg!

\endverse
\beginverse
Lad det gro
Kan ik' holde det hem'ligt mer'
Lad det gro
Det' på hagen det hele sker
Lad dem bare glo

\endverse
\beginverse
Alt det din krop ik' må
Alt det jeg ik' kan få

\endverse
\beginverse
Men i mit ho'ed er en idé
Det bli'r fantastisk, prøv at se
En datalog kan jeg nu bli'
Endelig

\endverse
\beginverse
Lad det gro
Lad det gro
Lad dem bare glo

\endverse
\beginverse
Mit skæg det snor sig som en slange mod mit bryst
De tykke lokker hvirvler sig i smukke cirklers lyst
Ved Nørrebros Runddel ligger min lykkens borg
Det' Fest og Farver, der har reddet os fra sorg

\endverse
\beginverse
Lad det gro
Lad det gro
Lad dem bare glo
Uden skæg er en datalogs dage talt'

\endverse
\endsong



\beginsong{Statistik}%71
[sr={Melodi: Frozen -- Let It Go}
,
by={Matematikrevyen 2015}
,
cr={}]\hypertarget{Statistik}{}
\label{song71}

\beginverse
Jeg ved ikke helt, hvad jeg skal gøre nu
Ingen folk, der viser vej
Det er jo en helt ny verden. Mon jeg tør? Åh nej!
Jeg ønsker folk, de ikke ville dømme mig
Nu vil jeg stå fast og være sej
\endverse

\beginverse
Nu gi'r jeg slip, nu er jeg klar
Ta'r en chance, det slut med algebra!
Er modig nu, det skal nok gå
Nu kan jeg få
\endverse

\beginverse
Statistik, hvorfor ik'?
Jeg elsker det som det er
Statistik, ren logik
Jeg vil bare gerne ha' mer'
Mat'matik det er ikke for mig
så nu vælger jeg
at t-testning det er min rette vej
\endverse

\beginverse
Når jeg ta'r min beslutning, er der ingen vej tilbag'
Så er det slut med venner, fællesskab og fredagskag'
Det bli'r farvel til mat'matik
Og så goddag til statistik
Farvel til skør geometri
Jeg er fri!
\endverse

\beginverse
Statistik, ren logik
Jeg vælger, hvad der er bedst
Statistik, kun teknik
Nu bli'r mit liv en fest
Når jeg står fast på, hvem jeg er
Kan jeg klare mer'
\endverse

\beginverse
En estimator, CLT og Theta-hat,
Normalfordeling og et histogram er vores skat
Min tvivl er væk og nu står kandidaten klar
Vil aldrig vende rundt. Farvel til det der var!
\endverse

\beginverse
Statistik, ren logik
Jeg er klar til nye fag
Statistik, nul kritik
Det er fedt hver en'ste dag
Så kom med, stud'et bli'r en leg
Og du bliver sej,
For R-kodning det er den sande vej!
\endverse
\endsong



\beginsong{Eksamensangst}[sr={Melodi: Anders Matthesen - Du er min store kærlighed}
,
by={FysikRevy 2014}
,
cr={}]\hypertarget{Eksamensangst}{}
\label{song72}

\beginverse
Uret tikker, hjertet det slår,
jeg mærker pulsen inde i mit øre.
Jeg får det værre som tiden går,
men jeg har lovet at jeg nok forsøger.
Jeg har dumpet alle mine fag.
Jeg håber i dag er min lykkedag.
\endverse

\beginverse
Solen falder blidt på min kind,
og jeg mærker trykket i min blære.
Jeg kan ikke nå at lette det nu;
nu må det briste eller bære.
Døren knirker, jeg gi'r et lille hvin.
Jeg håber snart at turen er min.
\endverse

\beginverse
\emph{[Omkvæd:]
Bare jeg nu kan få det sagt,
det jeg har øvet og planlagt,
det jeg har terpet ud og ind,
eller går klappen ned igen?}
\endverse

\beginverse
Hånden ryster, halsen er tør.
Jeg sveder kraftigt under mine arme.
Jeg stirrer tomt lige ind i på den dør
der står imellem mig og helvedets varme.
De andre gør det gang på gang på gang,
for mig er det stadig ulidelig tvang.
\endverse

\beginverse
\emph{[Omkvæd]}
\endverse

\beginverse
Bare jeg ku forklare,
bare censor ku' prøv' at forstå,
at mundtlig eksamen ik' er sag'n når man bli'r så nervøs at man går i stå.
\endverse

\beginverse
\emph{[Omkvæd]}
\endverse

\beginverse
\emph{[Omkvæd (moduleret)]}
\endverse
\endsong



\beginsong{Jeg har set en rigtig .*}[sr={Melodi: Jeg har set en rigtig negermand}
,
by={DIKUrevy/Sommerfest 2002}
,
cr={}]\hypertarget{Jeg har set en rigtig .*}{}
\label{song73}

\beginverse
Jeg har set en rigtig fysiker.
Han havd’ skæg og hår så langt som en musiker.
Han var sikkert ikk’ en rigtig mand.
Han ligned’ en der sku’ ha buksevand.
Jeg spurgte ham: "Hvad er du da for en?
Hvorfor har du så tynde blege stankelben?
Kom du fra rummet? Hvem er mon din mor?"
Så lo han blot og sagde disse ord:
\emph{``Det er sandt, at I fandt en fysikmutant.
Jeg har tragten med og hjernen sat i pant.''}

\endverse
\beginverse
Jeg har set en rigtig smart jurist.
Gik i jakkesæt og spilled sej, det var sgu’ trist.
Det kan godt vær’ han får mange pi’r,
men hans studiejob det er at ta’ kopier.
Hans hobby er at mishandle små dyr.
Før han vil hjælpe nogen kræver han gebyr.
Men det er også klart han er sadist,
det skal man være for at bli’ jurist.
\emph{En jurist er så trist, han er AntiKrist.
Gider nogen hente Max von XORcist?}

\endverse
\beginverse
Jeg har set en rigtig humanist.
Han var knokleskæv og sikkert også kommunist.
Så jeg spurgte hvad hans pensum var,
og han kigged’ op og gav mig dette svar:
``Til sommer skal jeg op i Anders And,
et nummer af Fantomet og en Superman,
og hvis en dag jeg vågner af min døs,
så bli’r jeg færdig og er arbejdsløs.''
\emph{Har meldt pas! Tigger hash! Lever kun på nas!
Når vi styrer verden får I alle gas!}

\endverse
\beginverse
Jeg har set en rigtig biolog,
og hun rulled’ sig i mudder’t som en anden so.
Hendes hår var fyldt med andemad.
Hun var måske køn hvis bar’ hun fik et bad.
Jeg spurgte: "Er du kvinde eller mand?
Hvorfor har du det mudder nede i din spand?
Har du en kær’ste, kan jeg bli’ din fyr?"
"Nej tak," sa’ hun, "jeg boller kun med dyr."
\emph{Biolog, er en so, sku’ man ellers tro.
Hvis du ellers får hend’ ren, så er hun go’.}

\endverse
\beginverse
Jeg er go’, for jeg er datalog.
Jeg har jakkesæt \& slips \& nye gummisko.
Mine bumser fik lidt kirurgi,
og nu er det mig som pigerne kan li’.
I Rungsted lever jeg i sus og dus
med Porsche, fast forbindelse og kæmpe hus.
Jeg lever højt på Internet-trafik,
og tjener fedt på andre fjolsers klik.
\emph{Jeg var fed, jeg var led, på succes jeg red
indtil sidste år da aktierne gik ned.}

\endverse
\endsong



\beginsong{På mat'matik}[sr={Melodi: Den Lille Havfrue -- Havet er skønt}
,
by={Matematikrevyen 2007}
,
cr={}]\hypertarget{På mat'matik}{}
\label{song74}

\beginverse
Ja, kurser er altid nemmer’
på et andet fakultet;
men russer, det her I glemmer
at KUA er ildeset.
Så selv hvis en dag du dumper
et fag og det er lidt trist;
så tænk på de andre tumper
der ikke får spillet whist!
\endverse

\beginverse
På mat’matik! På mat’matik!
Det’ så moderne, at gi’ sin hjerne
lidt gymnastik.
Vi har Kom-An og Algebra.
Vi lægger til og trækker fra.
Alle de regner, indtil de segner
på mat’matik!
\endverse

\beginverse
Så hvis du vil være doktor
og skære i døde folk,
så gi’r vi et samlet: ”Fuck dig!”
og din nekrofile dolk!
Nogen vil sygepleje
og gi’ andres tarm et skyl;
men her har vi drukspilslege
og tragter en masse øl!
\endverse

\beginverse
På mat’matik! På mat’matik!
Ingen, der koder, ingen, der roder
som på fysik.
For deres strenge teori
den er der ingen mening i.
Gå til ekstremer, vis teoremer,
på mat’matik! (På mat’matik)
På mat’matik (På mat’matik)
Her kan du lære at integrere
uden panik.
For vores fag det er så smukt
selv det der aldrig bliver brugt.
Vi vil studere mere og mere
på mat’matik!
\endverse

\beginverse
Polit - det er skidt.
For dansk eller spansk
Det’ plat - de er skvat
Og tysk noget pjat
Latin - er til grin
Fysik sutter *biip*
Og idræt - det rene råd
Kemi - nul magi
Dat-hum - er du dum?
Foragt for stat-akt
Jurist - det er trist
En flirt med en nørd
En fejl uden gejl
Og nano er virk’lig skod!
\endverse

\beginverse
Men mat’matik! Ja mat’matik!
Ingen kan bringe dig en så ringe
pædagogik!
Husk at din forelæsers ord
de betyder ikke spor.
Når han dig tugter, brug din instruktor
på mat’matik!
Øller i byen og til revyen
på mat’matik!
Drop nu kemien, biologien
økonomien, geografien
(sociologien, biokemien)
(teologien, filosofien)
(etnologien, psykologien)
(datalogien, fysiologien.)
Find drømmepigen
og mærk magien
på mat’matik!
\endverse
\endsong



\beginsong{Der er et ølrigt land}[sr={Melodi: Der er et yndigt land}
,
by={}
,
cr={}]\hypertarget{Der er et ølrigt land}{}
\label{song75}

\beginverse
Der er et ølrigt land,
det står med nød og næppe
blandt alt det pokkers vand
blandt alt det pokkers vand.
Det bugter sig i bar og kro.
Det hedder gamle Danmark,
og her er øllen go'
ja, hver en øl er go'.

\endverse
\beginverse
Her drak i fordums tid
hver tillakkede kæmper
sin mjød af fad med flid
sin mjød af fad med flid.
Så prøved' han, ej uden mén
at finde sine bene,
men faldt ved hver en sten
ja, hver en bautasten.

\endverse
\beginverse
Den øl endnu er skøn,
og gid den aldrig vælter.
Lad baj'ren stå så grøn
lad baj'ren stå så grøn.
De ædle sorters skønne øer
med sutter, sulde svende
og svimle danske møer
ja, svimle danske møer.

\endverse
\beginverse
Hil druk og fædreland.
Hil hver en kølig bajer.
Vi drikker dem vi kan
vi drikker dem vi kan.
Vort gamle Danmark -- SKÅL! -- bestå
så længe øllet skummer,
og næsen den bli'r blå
med røde prikker på.

\endverse
\endsong



\beginsong{Kursustilmelding}
[sr={Melodi: Christian Brøns -- Du kan gøre hvad du vil}
,
by={Matematikrevyen 2015}
,
cr={}]\hypertarget{Kursustilmelding}{}
\label{song76}

\beginverse
Havde du forventet
At du ku' vente
Med kursustilmelding?
Nu' det for sent
Du var forvænt med
Masser af frihed
Nu er den tag't fra dig
Væn dig til det
\endverse

\beginverse
\emph{Du kan vælge hvad du vil
Så længe du gør det på
Normeret tid, min ven
Du kan vælge hvad du vil
Så længe du gør det for
Reformen}
\endverse

\beginverse
Min forelæser sagde
Har du ambitioner?
Så tag dette kursus
I næste blok
Men det er umuligt
Jeg valgte jo kurser
For seks måneder siden
Ka' ik' laves om
\endverse

\beginverse
Nu er jeg bundet
På hænder og fødder
Bli'r gelejdet igennem
Alle fem år
Kan ikke læng're
Se nogen udvej
Men heldigvis slipper
Jeg for at tænke selv
\endverse

\beginverse
\emph{Du kan vælge hvad du vil
Så længe du gør det på
Normeret tid, min ven
Du kan vælge hvad du vil
Så længe du gør det for
Reformen}
\endverse

\beginverse
\emph{\lrep Du kan vælge hvad du vil
Så længe du gør det på
Normeret tid, min ven\rrep
Du kan vælge hvad du vil
Så længe du gør det for
Reformen}
\endverse

\beginverse
Havde du forventet
At du ku' vente
Med kursustilmelding?
Nu' det for sent
Du var forvænt med
Masser af frihed
Nu er den tag't fra dig
Væn dig til det
\endverse

\beginverse
Min forelæser sagde
Har du ambitioner?
Så tag dette kursus
I næste blok
Men det er umuligt
Jeg valgte jo kurser
For seks måneder siden
Ka' ik' laves om
\endverse
\endsong



\beginsong{Kvantepartiklen}[sr={Melodi: ShuBiDua - Den røde tråd}
,
by={FysikRevy 2014}
,
cr={}]\hypertarget{Kvantepartiklen}{}
\label{song77}

\beginverse
Hvor mon er før man bli'r set?
I Timbuktu, måske Tibet?
Fordeling af sandsynlighed?
Man må for fa'n have været et sted!

\endverse
\beginverse
I dobbeltspalten var kun mig.
Hvordan sku' jeg vælge vej?
Og jeg indså til mit held
inteferensen med mig selv.

\endverse
\beginverse
Potentialesymmetri
kan være ulig' eller lig'.
Der findes intet bedre sæt,
du ved min basis er komplet.

\endverse
\beginverse
%[Omkvæd:]
\emph{Helt nede
Livslede
Hvorfor kan jeg ej kendes?
Jeg længes efter vished, efter sikkerhed.}

\endverse
\beginverse
Endli' fik jeg position,
blev til en deltafunktion.
Operator'n gav et sted.
Nu har jeg endelig fået fred.

\endverse
\beginverse
\emph{Helt nede
Livslede
Min position forsvinder.
Kun minder om at være endegyldigt kendt.}

\endverse
\beginverse
Hvad mon man er før man bli'r set?
Partikel-bølge dualitet?
Bølgefunktion, en fermion?
Blot matematisk abstraktion.

\endverse
\endsong



\beginsong{Ode til bacon}%78
[sr={Melodi: Backstreet boys -- I want it that way}
,
by={Rasmus Fruergaard-Pedersen, Julerevy 2005}
,
cr={}]\hypertarget{Ode til bacon}{}
\label{song78}

\beginverse
Gik hen i Fø\TeX.
Det var en refleks.
Det var med en ræson.
Jeg sku’ ha’ bacon.

\endverse
\beginverse
Jeg så i disken,
helt tom. Jeg følte
på min vom – var som beton.
Den skal ha’ bacon.

\endverse
\beginverse
\emph{Hvorfor er
der tomt i køledisken?
Hvorfor er
der ingen lækkerbidsken?
Og hvor er
den gris der’ skåret i facon?
Jeg vil ha’ bacon.}

\endverse
\beginverse
Jeg så i vantro.
(Vægt)-vogterne de stod og lo.
Kom nu herhen smag vor karton.
Nej, jeg vil ha’ bacon.

\endverse
\beginverse
\emph{Hvorfor er
der tomt i køledisken?
Hvorfor er
der ingen lækkerbidsken?
Og hvor er
den gris der’ skåret i facon?
Jeg vil ha’ bacon.}

\endverse
\beginverse
Smagen kan give dig lysten til livet,
selvom du er vegetar, Aaah!
Du blir hvad du spiser, så blir jeg en gris, ja.
Det’ hel’re end hundred’ år!

\endverse
\beginverse
Stegt let med smør til.
Svøbt om pølser på grill.
Stegt sprødt,
stegt sprødt, stegt sprødt, stegt sprødt\ldots
Svøb flæsk i bacon!

\endverse
\beginverse
Stegt i smør eller helt naturel.
Føler mig så høj og kulturel.
Jeg ku’ synge ud fra en balkon:
``Jeg vil ha’ bacon.''

\endverse
\beginverse
\emph{Smag på lidt
bacon svøbt om dejlig mørksej.
Smag på lidt
bacon fyldt i julepostej.
Smag på lidt
bacon strøget let med stærk dijon, dijon.
Smag på lidt bacon.}

\endverse
\beginverse
\emph{Giv mig en
bacon-krydder-frikadelle.
Jeg har lyst
til en saltet stegt fedtcelle.
Hvem laver
mig en saltet flæskesværs-bonbon?
Jeg hylder bacon.
For jeg vil ha’ baco-o-o-n-n-n-n.}

\endverse
\endsong



\beginsong{Server'n er Crashed}%79
[sr={Melodi:Backstreet Boys - I Want It That Way}
,
by={DIKUrevy 2011}
,
cr={}]\hypertarget{Server'n er Crashed}{}
\label{song79}

\beginverse
Yeah!
Jeg bli'r
helt ked, men
Nu er
den nede
Den var belastet
Nu er den crashed.

\endverse
\beginverse
Min prompt
den sejler.
System-
et fejler
Der er sgu knas med
den server, der crashed.

\endverse
\beginverse
For åh nej.
Vi sku' ha købt et nyt raid
For åh nej
Jeg var sgu' ikke beredt.
før vi så
undtagelsen den kasted.
Server'n er crashed

\endverse
\beginverse
Men kan
den tvinges
til at
ku' pinges
Nej, den er - helt trashed
ja, server'n er crashed.

\endverse
\beginverse
For åh nej ...

\endverse
\beginverse
Tænkte, jeg fikser det bare i mor'n
da den stod der og kasted' fejl
yeah
Jeg burde ha' tjekket, men surfede porn
Nu er den sgu' gåe't sin vej.

\endverse
\beginverse
Det' ikk'
så skid' rart
Jeg ta'r
en genstart
Men nej
åh nej
åh nej
åh nej
Det fucking neder'n!

\endverse
\beginverse
Den fejler nu ved startup
Vi har sgu ingen backup
Disken er sikkert kvæstet
Server'n er crashed

\endverse
\beginverse
For åh nej ...
For åh nej ...

\endverse
\beginverse
Server'n er crashed

\endverse
\endsong



\beginsong{Koder Kapitalen}[sr={Melodi:Kim Hot - Smoker Industrien}
,
by={DIKUrevy 2013}
,
cr={}]\hypertarget{Koder Kapitalen}{}
\label{song80}

\beginverse
For 20 år siden der vidste folk ik'
at datalogi det ku' blive så fedt!
Der var en mand, der vidste det lidt
men han koded' slam og fortalt' det ik
men den mand han ku' se, at der kom et sprog
hvor alle der koded' ku' ligne et fjog!
Det fucking køter sprog det kørt' i en VM
og med \texttt{static void main} bli'r det hele så nemt!

\endverse
\beginverse
Det var 20 år siden og hans navn var Karl Koder
koded' slam for kapitalen, havde syge frynsegoder
Bawlet i Eclipse og spurgt': "Har du et problem?"
Så skal Karl nok kom' og kode dit system.

\endverse
\beginverse
Der var en konkurrence hvor han ik' sagde slut,
for Karl han sgu battle mod motherfucking Knuth!
Hans første linjer kode de var ik' særligt dope,
men resten de var bomben, for de brugte nemlig SOAP!

\endverse
\beginverse
Karl vandt og ham Knuth fik piveren på!
Tog tilbage til Stanford den BØSSEKAL!
Karl han sagde til den anden: "Pas godt på køter Java",
for det' fremtidssproget, der bliver fuldstændig fucked,
vi skal kode det hele! Forstår du det?
Og det skal kodes i Java, der må ik' vær' noget C!
Så fuck købsaftalen, tør røv med manualen
navnet er Køter Karl, og vi koder kapitalen!

\endverse
\beginverse
Koder ik i C,
heller ej VB,
eller det der lodne PHP
han koder Java,
kongen af Java,
han koder altid,
alting i Java!

\endverse
\beginverse
Koder ik i C,
heller ej VB,
eller det der lodne PHP
han koder Java,
kongen af Java,
kongen af Java,
kongen af Java!

\endverse
\beginverse
Nu står jeg her igen! Mand mod mand
Parat til at kod' hvad du ik' kan!
For datalogi det er blevet så fedt,
man kan lave en service på internet!
Og når den bli'r sat op, så kører den på Tomcat
og så ta'r Karl hjem, og holder sig en sabbat,
og så si'r han til sin chef:
"Fuck din kode! Det er UML din BØSSEKAL!"
Det skal kodes i Java! Det skal kodes med SOAP
I en managed verden uden at tænke på scope.
Hvor man elsker Swing, og bli'r fuldstændig fucked,
Så længe man koder Java, så går det godt!
Så alle de svanser, som der koder det i VB
Det skal nok blive godt med nog't Java på jeres CV
Så tag hjem fra regn'centralen og kod kapitalen!
Fræse lortet ned med SOAP, regulere hele dagen!

\endverse
\beginverse
Koder ik i C ...
Koder ik i C ...

\endverse
\beginverse
Når nu dikufanter de vil sælge kod' til nogen,
kan de lære et par hacks inde på Kongens kodeskole
Det første øvel'shold, det hedder SOAP-vold
Man bli'r låst inde i Hel, og skal håndkod' XML!
Den anden disciplin, den er ik' så nem at gætte,
man bli'r tævet med klassestil til man ka' Javaetikette.
Den tredje disciplin, det' alt for fedt
Du skal kode noget Java for HI-PER-FIT.
Man sletter GHC, starter Java-IDE
Fucker deres lambda, det en syg idé
Man lærer de svanser at hyle som en kat
når man breaker deres build hver eneste nat
Den allersidste time er lavningsdisciplin,
for når du koder Java, så bli'r damerne til dine,
Så fuck kvaliteten og skid på omtalen
Foran terminalen kan du kode kapitalen!

\endverse
\beginverse
Koder ik i C ...
Koder ik i C ...
Koder ik i C ...
Koder ik i C ...

\endverse
\beginverse
Koder Karl...

\endverse
\endsong



\beginsong{Minken Mink}[sr={Melodi:Bjørnen Bjørn er en bjørn}
,
by={Computer Science Camp 2015}
,
cr={}]\hypertarget{Minken Mink}{}
\label{song81}

\beginverse
Minken Mink er en mink.
Minken Mink er en mink.
Og den kan li' de fleste, den er så flink,
For Minken Mink er en mink

\endverse
\beginverse
Nogen tror at Minken er en ræv.
Nogen tror at Minken er en husmår.
Nogen tror at Minken er en kat.
Men Minken Mink er en mink.
Ja Minken Mink er en mink!

\endverse
\beginverse
Minken Mink er en mink.
Minken Mink er en mink.
Og den kan li' arrangører for de er pink
For Minken Mink er en mink

\endverse
\beginverse
Nogen tror at Minken er en ræv.
Nogen tror at Minken er en husmår.
Nogen tror at Minken er en kat.
Men Minken Mink er en mink.
Ja Minken Mink er en mink!

\endverse
\endsong



\beginsong{Drop Nu Ud}[sr={Melodi:The Lion King - Be Prepared}
,
by={FysikRevy 2011}
,
cr={}]\hypertarget{Drop Nu Ud}{}
\label{song82}

\beginverse
I ved at jeg er nederdrægtig
Jeg er ond som en anden dekan
Ja, min sjæl er der noget fordækt i
Men hør nu min udsøgte plan

\endverse
\beginverse
Med mindre du er overlegen
Hvis ikke du er et geni
Er det bedst at du skaffes af vejen
Og se her har jeg en strategi:

\endverse
\beginverse
Drop ud - hvis du ikke kan kvante
Hvis din IQ er under en-treds
Jeg udrydder tumper, så studiet skrumper
(Og hvad hvis jeg dumper?)
Jeg flår dig i stumper!
Hvis du er uværdig og aldrig bli'r færdig
Vil jeg knuse dit mod fra dag et
Og kun du, lille rus, står for skud
Drop nu ud!

\endverse
\beginverse
(Hah! Hvorfor skulle vi droppe ud? Vi er da helt vildt skarpe)
I vil måske forsøge jer med et eksamensspørgsmål her og nu?
(Bare kom an Mogens!)
Nuvel. Udregn energierne for en partikel i en brønd.
(Hurra! Vi bruger bare Newtons love!)
(Jaaa! Newton, Newton, la lalalalala!)
IDIOTER! I kan ikke løse det uden Schrödingerligningen.
(Jamen - det er jo vanvid?)
Vanvid?
DET. ER. KVANT!!!

\endverse
\beginverse
(Argh! Skynd jer at lære kvant! Ellers får vi aldrig nogle ECTS!)
(ECTS!)
(ECTS!!)

\endverse
\beginverse
(Nu terper vi kvant hele dagen)
(Og natten går med QCD.)

\endverse
\beginverse
Og sku I få lyst til at klage
er der intet at gøre ved det.

\endverse
\beginverse
For søger du inspirationen
så siger jeg 'the cake is a lie'
Men mangler du motivationen,
Så husk blot at 'terpen macht frei!'

\endverse
\beginverse
Så drop ud - hvis du gerne vil skånes
Drop ud, hvis du har livet kært.
For her ruller hoveder - bag tonede ruder (Terp nu kvant, terp nu kvant)
Du trygler 'vær nådig' - der runger af gråd i (Terp nu kvant, terp nu)
eksamenslokalet, din død bliver forhalet
og spør' du til din karakter
vil jeg ridse 00 i din hud
-Drop nu ud!

\endverse
\beginverse
-Så giv op, fald på knæ, rus, og tud
-Drop nu ud!
(Mwuhuhahahah!)

\endverse
\endsong



\beginsong{Primtal}[sr={Melodi: Robbie Williams - Candy}
,
by={Matematikrevyen 2014}
,
cr={}]\hypertarget{Primtal}{}
\label{song83}

\beginverse
Lad os definere
de tal der fascinerer.
For primtal skal der gælde
faktorer trivielle.
Elegant og simpelt, det kan man let forstå.
Men trods det er det fulde billed’
umuligt at opnå.
Men hør nu:
\endverse

\beginverse
Euler og Euklid de viste,
blandt primtal er der slet intet sidste.
Givet $n$ kan primtal skrives
helt entydigt, vi har uniqueness.
\endverse

\beginverse
Skriv nu $p$ og $q$.
Der’ så meg’t der ikk’ er vist endnu.
Uden dem går teori itu,
for alt er smukt ved primtal.
Skriv nu $p$ og $q$.
De er overalt, det vides jo.
Også selvom det er svært at tro,
for alt er smukt ved primtal.
\endverse

\beginverse
Deres heltalsringe
dem kan man let frembringe.
For vilkårligt $p$
er der meget at indse.
Man kan se tendenser
og smukke kongruenser.
Perfekt at anerkende,
forbundet med Mersenne.
\endverse

\beginverse
Goldbach havde en formodning.
Og kryptologer brug’r dem i kodning.
Dankortkøb forløber sikkert fordi
vi slipper primtal fri.
Hvad gjord’ vi uden dem?
Så kom nu!!
\endverse

\beginverse
Skriv nu $p$ og $q$\ldots
\endverse

\beginverse
$\pi(x)$ er asymptotisk.
Distribueringen dog ej logisk.
Sylows sætning hjælper os til at se
grupper af orden $p$.
%\vleftofline{9x }
Hvad gjord’ vi uden dem?\rep{9}
%Hvad gjord’ vi uden dem?
%Hvad gjord’ vi uden dem?
%Hvad gjord’ vi uden dem?
%Hvad gjord’ vi uden dem?
%Hvad gjord’ vi uden dem?
%Hvad gjord’ vi uden dem?
%Hvad gjord’ vi uden dem?
%Hvad gjord’ vi uden dem?
\endverse

\beginverse
Skriv nu $p$ og $q$\ldots
\endverse

\beginverse
Skriv nu $p$ og $q$\ldots
\endverse
\endsong



\beginsong{Kringsatt av fiender}[sr={}
,
by={Tekst: Nordahl Grieg -- Melodi: Otto Mortensen}
,
cr={}]\hypertarget{Kringsatt av fiender}{}
\label{song84}

\beginverse
Kringsatt av fiender, gå
inn i din tid!
Under en blodig storm,
vi deg til strid!
 Kanskje du spør i angst,
udekket, åpen:
 ``Hva skal jeg kjempe med,
hva er mitt våpen?''
\endverse

\beginverse

Her er ditt vern mot vold,
her er ditt sverd:
 troen på livet vårt,
menneskets verd.
 For all vår fremtids skyld,
søk det og dyrk det,
 dø om du må, men
øk det og styrk det!
\endverse

\beginverse

Stilt går granatenes
glidende bånd.
 Stans deres drift mot død,
stans dem med ånd.
 Krig er forakt for liv,
fred er å skape.
 Kast dine krefter inn,
døden skal tape.
\endverse

\beginverse

Elsk - og berik med drøm -
alt stort som var.
 Gå mot det ukjente,
fravrist det svar!
 Ubygde kraftverker,
ukjente stjerner –
 skap dem, med skånet livs
dristige hjerner.
\endverse

\beginverse
Edelt er mennesket,
jorden er rik.
 Finnes her nød og sult,
skyldes det svik.
 Knus det! I livets navn
skal urett falle.
 Solskinn og brød og ånd
eies av alle.
\endverse

\beginverse
Da synker våpnene
maktesløs ned!
 Skaper vi menneskeverd,
skaper vi fred.
 Den som med høyre arm
bærer en byrde,
 dyr og umistelig,
kan ikke myrde.
\endverse

\beginverse

Dette er løftet vårt
fra bror til bror:
 Vi vil bli gode mot
menneskenes jord.
 Vi vil ta vare på
skjønnheten, varmen –
 som om vi bar et barn
varsomt på armen.
\endverse
\endsong



\beginsong{Smukke ting ved mat'matik}[sr={Melodi: Pocahontas - Colours of the Wind}
,
by={Matematikrevyen 2008}
,
cr={}]\hypertarget{Smukke ting ved mat'matik}{}
\label{song85}

\beginverse
I synes, at vi ik' gavner meget
vor's fag er teoretisk...
\hspace{5em}og det er også sandt
Men det I ik' forstår
er den følelse vi får
når vi i teori kan løse alt
løse alt...

\endverse
\beginverse
Vi regner ikke for at redde verden
Vi regner ikke for at blive rig'
Vi elsker bare alt ved mat'matiken
Det' den verden, som vi godt vil være i

\endverse
\beginverse
Vi accepterer andre læser andet
så vi kan læse vores mat'matik
og opdage det smukkeste i verden
brug' den chance, som I andre aldrig fik

\endverse
\beginverse
Har I aldrig ladet log gå mod uendelig
og set hvad vej at den nu egentlig gik
Kan I ikke finde charmen ved en ligning
Kan I se de smukke ting ved mat'matik

\endverse
\beginverse
Kom lad os integrere alle udtryk
og undersøg' 1/x i nul
og undersøge alt med asymptoter
finde ud af, hvad der er i dette hul

\endverse
\beginverse
Vor's sprog det er præcist og meget tyd'ligt
og misforståelser de sker ej her
og det er det, der gør, at mat'matikken
er et fag, der kan nå så meget mer'

\endverse
\beginverse
Gammafordelingen
den er måske svær...men den er din ven

\endverse
\beginverse
Vil I aldrig lade log gå mod uendelig
forsøg selv om det bryder din etik
Kan I ikke finde charmen ved en ligning
Kan I se de smukke ting ved mat'matik
læs nu blot på jeres fag
men det sker nu nok en dag
at I ser\ldots de smukke ting ved\ldots mat'matik\ldots

\endverse
\endsong


\beginsong{Induktionssang}[sr={Melodi: Kim Larsen og Bellami - Familien skal i skove}
,
by={Matematikrevyen 2015}
,
cr={}]\hypertarget{Induktionssang}{}
\label{song86}

\beginverse
Det er efterår, jeg skal til Dis-eksamen
Jeg har læst og læst, helt ifølge planen
Snupper min cykel – af sted! Nu' der ingen vej tilbage
Dugen, den er grøn, censor, hun nynner
Hilser med hånden:
\emph{Inge-Lise Brønner}
Til tavlen går jeg hen, det' jo dér, hvor det hele, det begynder

\endverse
\beginverse
Eksaminator spørger ind til induktion
Så jeg si'r: ``Peanos ak-si-om''
\endverse

\beginverse
Det har jeg brugt så mange gange før
Og nu gør jeg det igen
\emph{Det har han brugt så mange gange før
Og nu gør han det igen}

\endverse
\beginverse
Gælder det for $n$, så gør det ogs' for resten
\emph{Ingen induktionsstart -- hør nu her, nu stopper festen}
Hvorfor kan jeg ikk' bar' la' $n$'et flyve ud i det uendelig'?
\emph{Åh! Dumme dreng!}

\endverse
\beginverse
Antag $n$ er 1, for da må det gælde
Jo, det er da rigtigt, nu skal du bare tælle
Det kan jeg da sagtens: 1, 2, 3, 4, 5 og ælle, bælle
\emph{Prøv nu at hold' fokus, og sig os så, min ven
Hvordan går du fra $n$ til $n + 1$?}

\endverse
\beginverse
Det har jo gjaldt for så mange andre $n$
Så det gør det nok igen
\emph{Det har jo gjaldt for så mange andre $n$
Så det gør det nok igen}

\endverse
\beginverse
Det er blevet sent, de skal votere
Så mindes jeg med gru mine dumpekarakterer
Håber sådan dette bli'r 02, eller måske endda 4
Døren, den er lukket, gangene er øde
Var det tilstræk'ligt, får jeg mon det røde?
Tiden er gået, nu sker det, tilbage i lokalet
Censor, hun siger:
\emph{Ja, nu må du gerne gå
Men vi ka' jo ikk' la' dig bestå}

\endverse
\beginverse
Det har de sagt så mange gange før
Derfor gør de det igen
\emph{Det har vi sagt så mange gange før
Derfor gør vi det igen}
Det har de sagt så mange gange før
Derfor gør de det igen
\emph{Det har vi sagt så mange gange før}
Derfor gør de/\emph{vi} det igen


\endverse
\endsong



\beginsong{Everything is Gaussian!!!}[sr={Melodi: Everything is Awesome!!!}
,
by={ FysikRevy 2016}
,
cr={}]\hypertarget{Everything is Gaussian}{}
\label{song87}
\beginverse
		Everything is gaussian
		Everything you measure has a width and a mean
		Everything is gaussian: central limit theorem
\endverse
\beginverse	
		Nobody likes the poisson distribution
		Binomial is worse and we need a solution, what's the solution?
		N is big, P is small, you will see, your data will set you free
\endverse
\beginverse	
		Everything is gaussian
		Everything you measure has a width and a mean
		Everything is gaussian: central limit theorem
\endverse
\beginverse	
		(Wooo)
		3, 2, 1. Go
\endverse
\beginverse	
		Barlow and Taylor, they're just awesome
		Stat. made easy 'cause everything is Gaussian!
		Take the limit, N to infinity
		Convergence is guaranteed thanks to entropy
\endverse
\beginverse	
		I drew some awesome super pretty Gaussians
		Put my data into histograms
		while I listen to my favorit' jams
		Fittin' all the functions, everything is gaussian
		Changed the axes, got a new view
		It's gaussian to fit, and it's gaussian to do (It's gaussian to do)
\endverse
\beginverse	
		Nobody likes the poisson distribution
		Binomial is worse and we need a solution, what's solution?
		N is big, P is small, you will see, your data will set you free
\endverse
\beginverse	
		Everything is gaussian
		Everything you measure has a width and a mean
		Everything is gaussian: central limit theorem
\endverse
\beginverse	
		Bell curves, height of kings
		We just named two gaussian things
		Women's shoe size', table lengths
		You know what's gaussian? EVERYTHING!
\endverse
\beginverse	
		IQ tests, velocities,
		Number of people's allergies
		Wolfowitz' and Chi square test
		gasssian items are the best!
\endverse
\beginverse	
		Bells, Hills, Mills,
		They're gaussian
		Socks, Clocks and Rocks
		They're gaussian!
		Per Hede-gård
		He's gaussian!
		Everything you see, or think, or say
		Is gaussian!
\endverse
\beginverse	
		Everything is gaussian
		Everything you meassure has a width and a mean
		Everything is gaussian: central limit theorem
\endverse
\beginverse	
		Everything is gaussian
		Everything you meassure has a width and a mean
		Everything is gaussian: central limit theorem
\endverse
\endsong




\beginsong{Din TBF}[sr={Melodi: Village People - YMCA}
,
by={ UNFRevy 2016}
,
cr={}]\hypertarget{Din TBF}{}
\label{song88}
\beginverse
		Du der - du skal være arrangør
		Jeg sagde du der - hvis du ellers da tør
		Jeg sagde du der - hvis du vil være med
		til at gør' en kæmpe forskel
\endverse
\beginverse	
		Du der, hvis du skriver dit navn
		på den her er-klæring bliver dit savn
		for naturvi-denskab fuldstændig væk
		og du bliver en af os nu
\endverse
\beginverse	
		For du skal bare ha' en TBF
		Så bliver du en del af UNF
\endverse
\beginverse	
		Så kan du logge ind på vores eget system
		Du kan rode rundt i vores ting
\endverse
\beginverse	
		Ja, du skal bare ha' en TBF
		Så bliver du en del af UNF
\endverse
\beginverse	
		Men pas på hvad du si'r, hvis du let bliver genert
		For du bliver jo nok citeret
\endverse
\beginverse	
		Du der, har du fanget det nu?
		Jeg siger du der, er du klar til at brug'
		din tid på at sætte deadlines til folk
		som alligevel bryder dem, og
\endverse
\beginverse	
		har du nu helt styr på hvordan
		du kontakter forelæser så han
		siger ja tak så'n at vi kan få fyldt
		hele vores sæsonprogram
\endverse
\beginverse	
		For du har lige fået din TBF
		Så nu er du en del af UNF
\endverse
\beginverse	
		Referater er der, formularer med mere,
		og din pizzafaktor er her
\endverse
\beginverse	
		For du har lige fået din TBF
		Så nu er du en del af UNF
\endverse
\beginverse	
		Hvis du får brug for hjælp, eller bliver i tvivl
		kan en QuickPoll redde dit liv
\endverse
\beginverse	
		For du har lige fået din TBF
		Så nu er du en del af UNF
\endverse
\endsong


\beginsong{Et helt nyt liv}[sr={Melodi: Disney - A Whole New World}
,
by={ UNF Revy 2016}
,
cr={}]\hypertarget{Et helt nyt liv}{}
\label{song89}
\beginverse
        Jeg kan vise dig alt
		Denne strålende verden
		Hør, frivillig, hvornår har hjertet
		Sidst bestemt for dig
\endverse
\beginverse
        Jeg kan åbne dit blik
		mod naturvidenskaben
		Ja, du fatter det næppe
		Alt det som vil vise sig
\endverse
\beginverse
		Et helt nyt liv
		Et nyt fantastisk perspektiv
		Ingen der siger nej
		Og standser dig
		Og siger at vi kun drømmer
\endverse
\beginverse
		Et helt nyt liv
		En verden som er ny for mig
		I sort og hvid og blå, kan jeg forstå
		At viden kan formidles her, med dig
		(At viden kan formidles her, med dig)
\endverse
\beginverse
		Helt utrolige camps
		Ingen søvn og kun kaffe
		Venskaber kan jeg skaffe
		Mens jeg laver camps med dig
\endverse
\beginverse
		Et helt nyt liv
		(Tag en kuglepen)
		Og hundred' tusind ting at gøre
		(Bare vent der er mere)
		Jeg er en arrangør
		Jeg DK kører
		At vende om, er ingen god ide,
\endverse
\beginverse
        Mailen lukker sig op
		(Hvor det utroligste kan ske)
		Hvert projekt fascinerer
\endverse
\beginverse
		Jeg alle foredrag så
		Der nok at nå
		Lad mig få et helt nyt liv med dig
\endverse
\beginverse
		Et helt nyt liv
		Med gratis mad
		En helt ny vej
		Som gør mig glad
		Så tag med mig
		Til landsmøde
		For dig og mig
\endverse
\endsong


\beginsong{Hvad må man (som arrangør)?}[sr={Melodi: Omsen og Momsen - Fy fy skamme}
,
by={ UNFRevy 2016}
,
cr={}]\hypertarget{Hvad maa man}{}
\label{song90}
\beginverse
		Må man flyv’ på første klasse?
		Må man spise mad på NOMA?
		Må man hæld’ så meget sprut på delta’r’
		at de går i koma?
\endverse
\beginverse			
		Må man gamble med vores penge?
		Må man drukne sig i vodka?
		Må man bruge Times New Roman
		når man laver sig en tryksag?
\endverse
\beginverse			
		Næ næ næ næ næ det må vi ikke,
		fy fy skamme skamme fy fy ah ah
		slemme slemme fy fy næ næ nix nix
		slut forbudt - men hva må vi så?
\endverse
\beginverse			
		Må man ændre vores logo?
		Må man slet’ vores database?
		Må man pille på en deltager
		og gå til næste fase?
\endverse
\beginverse			
		Må man bade under sprinkleren
		som er i lab’ratoriet?
		Må man sparke til en forelæser
		nede i auditoriet?
\endverse
\beginverse			
		Næ næ næ næ næ det må vi ikke,
		fy fy skamme skamme fy fy ah ah
		slemme slemme fy fy næ næ nix nix
		slut forbudt - men hva må vi så?
\endverse
\beginverse			
		Må man sove til et for’drag?
		Må man ryge pot på brandvagt?
		Må man rode på vores lager
		eller arrangere en andagt?
\endverse
\beginverse			
		Må man udsmide plakater
		eller grine af folk der bløder
		Må man be’ om hemmelig afstemning
		til alle vores møder
\endverse
\beginverse	
		Næ næ næ næ næ det må vi ikke,
		fy fy skamme skamme fy fy ah ah
		slemme slemme fy fy næ næ nix nix
		slut forbudt - men hva må vi så?
\endverse
\beginverse			
		Må man melde sig til opvask
		Må passe på sin kasse
		Må man komm’ med gode input
		når vi planlægger en masse
\endverse
\beginverse			
		Må man sige fra i tide
		Må man stille op til DK
		Må man finde sig en UNF’er
		som man kan læg’ an på?
\endverse
\beginverse			
		Ja ja ja ja ja det må vi gerne
		næ sikke fint fint ja ja meget det blir
		stort stort klap klap fint fint flot flot
		det var godt - nå det må vi godt
\endverse
\endsong


\beginsong{Jeg elsker science}[sr={Melodi: Gloria Gaynor - I Will Survive}
,
by={ UNFRevy 2016}
,
cr={}]\hypertarget{Jeg elsker science}{}
\label{song91}
\beginverse
Først så var jeg bange, jeg var skrækslagen
Jeg tænkte, det lød ligeså skrækkeligt som et stræklagen
Og jeg sad søvnløs hele natten, stirred’ tomt ind i min pejs
Før hed det naturvidenskab, nu skal man sige “science”
\endverse
\beginverse
Syntes det var dumt - noget rigtigt lort
Men nu der ingen vej tilbage, “science” er det nye sort.
Jeg skulle ha sagt noget eller gjort noget, eller gået i protest
Men sket er sket, og science er jo dansk når det er bedst
\endverse
\beginverse
Det’ ikk’ så slemt, og egentlig nemt
Syv bogstaver? Det enkelt og bekvemt
Hvorfor skal man slæbe rundt på gamle danske ord?
Alle snakker engelsk nu, ligegyldigt hvor de bor.
\endverse
\beginverse
Jeg elsker science, jeg elsker science!
Fysik, kemi, biologi
det er jo alt for nice.
Matematik, geologi
og medicinsk teknologi, jeg elsker science
Jeg elsker science!
\endverse
\endsong


\beginsong{Science Summer Lovin'}[sr={Melodi: Grease - Summer nights/Summer lovin'}
,
by={ UNFrevy 2014}
,
cr={}]\hypertarget{Science Summer Lovin}{}
\label{song92}
\beginverse
		Dreng: Science Camp - had me a blast
		Pige: Science Camp - happened so fast
		Dreng: I met a girl - crazy for meeee
		Pige: Met a boy - smart as can be
\endverse
\beginverse
		High school days
		drifting away
		To, uh oh, those science camps
\endverse
\beginverse
		Well-a, well-a, well-a, huh
\endverse
\beginverse
		Teach me more, teach me more
		Did you learn a lot?
		Teach me more, teach me more
		Is science that hot?
\endverse
\beginverse
	Dreng: She stood by me, up in the lab
	Pige: My reactions - they went so baad
	Dreng: I saved her life - before it explod'
	Pige: He looked good, in his lab coat
\endverse
\beginverse
	Summer science
	something's begun
	And oh, oh, I wanna learn more
\endverse
\beginverse
		Well-a, well-a, well-a, huh
\endverse
\beginverse
	Teach me more, teach me more
	Can you help me with math?
	Teach me more, teach me more
	Does it turn out in cash?
\endverse
\beginverse
	Dreng: We were solvin' matrices and more
	Pige: Studin' Fermat - my head got sore
	Dreng: We did math - out in the park
	Pige: We stayed up till it turned dark
\endverse
\beginverse
	Summer science
	made friends for life
	and oh, oh, I wanna learn more
\endverse
\beginverse
	Teach me more, teach me more
	Can we play your new game?
	Teach me more, teach me more
	Will it lead you to fame?
\endverse
\beginverse
	Pige: My computer, holding my hand
	Dreng: summer indoor - didn't get tanned
	Pige: I made a game - I'm just eighteen
	Dreng: Well, it's good - you know what I mean
\endverse
\beginverse
	Summer camps,
	Science is fun
	And oh, oh, I wanna learn more
\endverse
\beginverse
	Teach me more, teach me more
	Anything you don't know?
	Teach me more, teach me more
	Next year I wanna go!
\endverse
\beginverse
	Pige: Had a party, that's where it ends
	Dreng: No more science, with my new friends
	Pige: A week of camp, that was just "wauw"!
	Dreng: Wonder what she's studying now?
\endverse
\beginverse
	Summer science, blowing my mind
	But oh, I wanna learn moooore
\endverse
\beginverse
	Teach me more
	teach me mooooore.
\endverse
\endsong



\beginsong{ScienceCamps}[sr={Melodi: tema fra Pokémon}
,
by={ UNFRevy 2016}
,
cr={}]\hypertarget{ScienceCamps}{}
\label{song93}
\beginverse
		Jeg drager ud på livets vej
		Jeg har et enkelt mål
		ScienceCamps afholder jeg
		Min vilje er af stål!
\endverse
\beginverse
		Starter ud på Chemisty
		Og ta’r så til Fysik
		GDC og medico
		derefter matematik
\endverse
\beginverse
		ScienceCamps, jeg kan klare dem!
		Multitasking er mit lod
		4 camps -åhh jeg elsker det
		bare sommer’n bli’r ved og ved
\endverse
\beginverse
		4 camps, jeg ka’ klare det,
		det kræver bare koffein
		Jeg vil bare hjælpe til!
		ScienceCamps - kan jeg klare dem?
		Jeg kan klare dem!
\endverse
\beginverse
		Socialt ansvar på nanocamp
		Kasserer på SDC
		Fagligt team og så PR
		Koordinators vilje ske!
\endverse
\beginverse
		Første uge er ovre nu
		Ser frem til kriminal
		Astro, Show og CSC
		Det bliver fænomenalt!
\endverse
\beginverse
		13 camps, jeg skal klare det!
		måske det er lidt hårdt
		20 camps,ååh, Discovery!
		Jeg skal smadre den rekord!
\endverse
\beginverse
		40 camps, jeg har brug for hjælp!
		Jeg mangler stadig CSI!
		robodays og futureweek
		medicin, jeg skal bruge det!
\endverse
\beginverse
		Når presset det bliver alt for stort
		og jeg går psykisk ned
		Så tager jeg bare en ekstra camp
		for jeg bliver ved og ved!
\endverse
\beginverse
		Min psykolog, han siger stop!
		- min læge ligeså
		Men jeg mangler biotech
		Den må jeg lige på!
\endverse
\beginverse
		100 camps, alle skal jeg nå!
		Det er det jeg duer til
		1000 camps,åh medicinal
		Robot og ISSC
\endverse
\beginverse
		Alle camps, åh, jeg er lidt træt
		Men stemmerne bliver ved og ved
		Næsteår, der står jeg klar!
		ScienceCamps, jeg kan klare dem,
		kan jeg klare dem? ScienceCamps
\endverse
\endsong





\beginsong{Om lidt bli'r her stille}%86
[sr={}
,
by={Kim Larsen}
,
cr={}]\hypertarget{Om lidt bli'r her stille}{}
\label{song94}

\beginverse
Om lidt bli’r her stille.
Om lidt er det forbi.
Fik du set det du ville?
Fik du hørt din melodi?
\endverse

\beginverse
Forladt og alene
danser cirkus prinsessen rundt.
Går i stå --- på sin line
i et sanseløst sekund.
\endverse

\beginverse
\emph{Om lidt om lidt
er vi borte.
Vi ses måske igen.}
\endverse

\beginverse
Om lidt bli’r her stille.
Om lidt er det forbi.
Fik du set det du ville?
Fik du hørt din melodi?
\endverse

\beginverse
\lrep\emph{Om lidt, om lidt
er vi borte.
Vi ses måske igen.}\rrep
\endverse
\endsong



\end{songs}\makeatletter
\newcount\divFour
\divFour \c@page
\divide\divFour by 4
\multiply\divFour by 4
\advance\divFour by-\c@page
\edef\remainFour{\the\divFour}
\ifnum\remainFour = 0
  \clearpage\null\clearpage\null\clearpage\null
\else\ifnum\remainFour = -1
  \clearpage\null\clearpage\null
\else\ifnum\remainFour = -2
  \clearpage\null
\fi\fi\fi
\newgeometry{margin=2cm,top=2cm,bottom=2cm}
\setlength{\headwidth}{\textwidth}
\showindex[2]{Sange}{titleidx}
\end{document}
